\section{Further layer instantiations}\label{app:further-instantiations}

This appendix collects the analogies that accompany the laws of the catalog, beyond those kept beside each law. They are grouped by law, in the order of the chapters, and all tagged as analogies; every special case is kept in the body.

\subsection*{\Fref{F10} Quotientality}
\begingroup\small
\begin{description}
\item[\emph{Trubetzkoy phonemic principle (structural linguistics).}] \textit{Analogy.}
For a language \(L\), the phoneme inventory \(\mathcal P_L\) is a
partition of phonetic realisations preserving every
morpheme-distinguishing minimal-pair contrast.  Two phones that
produce a minimal pair belong to distinct phonemes; phones that never
do are grouped into one phoneme only when they are also phonetically
similar (English \emph{h} and the final nasal of \emph{sing} never
contrast yet are distinct phonemes), so non-contrast alone is not an
equivalence.  Every morpheme-discriminating phonological observable
factors through the phoneme map, and merging two phonemes that form
a minimal pair collapses a morpheme distinction \(L\) tracks
\citep{Trubetzkoy1939GrundzugePhonologie}.

\item[\emph{Categorical perception of voice onset time (cognitive
psychophysics).}] \textit{Analogy.}
For a synthesised voice-onset-time continuum spanning \(/{\!b}/\) to
\(/{\!p}/\), human listeners exhibit \emph{categorical perception}:
within-category {ABX} discrimination is weak, falling toward but
generally not to chance, while between-category discrimination peaks
at the boundary near \(+25\)~ms.  Liberman and colleagues first
showed this pattern on a place-of-articulation continuum
(\(/{\!b}/\)--\(/{\!d}/\)--\(/{\!g}/\)).  The phonetic-category
partition is the coarsest equivalence on the acoustic continuum
preserving the identification probe, which factors through the
phonetic-category map; the ABX probe factors through it up to the
residual within-category discrimination
\citep{LibermanHarrisHoffmanGriffith1957Discrimination}.
\end{description}
\endgroup

\subsection*{\Fref{F11} Adequacy / No-Overread}
\begingroup\small
\begin{description}
\item[\emph{Post-processing invariance of differential privacy
(cryptography / privacy).}] \textit{Analogy.}
For an \((\varepsilon,\delta)\)-differentially-private mechanism
\(M:\mathcal D\to\mathcal R\), the post-processing property of differential privacy (proved for \(\varepsilon\)-DP by Dwork--McSherry--Nissim--Smith, and holding equally for \((\varepsilon,\delta)\)-DP) states that for any (possibly
randomised) \(g:\mathcal R\to\mathcal W\), the composition
\(g\circ M\) is itself \((\varepsilon,\delta)\)-DP.  Every post-processing
of \(M\) is therefore again \((\varepsilon,\delta)\)-DP; a published quantity whose
neighboring-database distinguishability exceeds the
\(e^{\varepsilon}\!\cdot\!\!\ (\cdot)+\delta\) bound is the
overread failure --- evidence of access outside \(M\)'s
randomisation
\citep{DworkMcSherryNissimSmith2006CalibratingNoise}.

\item[\emph{Registered Reports / preregistered analyses (scientific
methodology).}] \textit{Analogy.}
In the Registered Reports format the access quotient is the
\emph{Stage~1 preregistered protocol together with the locked
dataset}: confirmatory claims must be computable from those two
inputs alone, and any two execution histories sharing them must
produce the same confirmatory verdict.  HARKing, p-hacking, and
outcome switching are overread failures: a published claim
presented as confirmatory but computed from post-hoc analytic
choices fails to factor through the declared protocol-plus-data
quotient
\citep{Chambers2013RegisteredReports}.

\item[\emph{Federal Rules of Evidence: admissibility (legal
procedure).}] \textit{Analogy.}
Under the U.S. Federal Rules of Evidence (Rules 401--403 on
relevance, Rule 802 on hearsay), a finding of fact must factor
through the admitted-evidence quotient of the case record: two
factually distinct records with identical admitted-evidence content
must yield the same verdict.  A verdict resting on inadmissible
evidence is reversible error unless harmless, instantiating the
\Fref{F11} overread failure
\citep{FRE2024FederalRulesEvidence}.

\item[\emph{Halpern--Moses common knowledge (distributed-systems
epistemic discipline).}] \textit{Analogy.}
For a process \(i\) in a distributed system with local-state map
\(r_i\) and induced indistinguishability \(\sim_i\), the Kripke
semantics for \(K_i\varphi\) says that the truth of
\(K_i\varphi\) at a global state factors through \(r_i\): two
global states with the same local state for \(i\) make
\(K_i\varphi\) hold or fail together.  Accepted knowledge claims
by process \(i\) are therefore exactly those determined by its
view; a claim that depends on global state beyond \(r_i\) is the
overread failure.  Common knowledge \(C_G\varphi\) for a group \(G\)
is the group-level case: its truth factors through the finest
partition coarser than every local view \(r_i\), \(i\in G\), so it
holds or fails together on each connected component of the union of
the relations \(\sim_i\)
\citep{HalpernMoses1990KnowledgeCommonKnowledge}.
\end{description}
\endgroup

\subsection*{\Fref{F12} No-Free-Distinction}
\begingroup\small
\begin{description}
\item[\emph{Wolpert no-free-lunch theorem (statistical learning).}] \textit{Analogy.}
For supervised learning on a finite input space, averaged uniformly
over all possible target functions, every learning algorithm
achieves the same expected off-training-set error.  Off-sample
generalization is not a distinction obtainable from training data
alone; it requires the declared apparatus extension of an
inductive bias --- a hypothesis class restriction, a prior, or a
regulariser --- without which Wolpert's averaging identity collapses
every learner to chance
\citep{Wolpert1996NoFreeLunch}.

\item[\emph{Wootters--Zurek no-cloning theorem (quantum mechanics).}] \textit{Analogy.}
No unitary on \(\mathcal H_S\otimes\mathcal H_A\) sends
\(|\psi\rangle|0\rangle\) to \(|\psi\rangle|\psi\rangle\) for every
unknown \(|\psi\rangle\); the linearity of unitary evolution
forbids perfect copying of an unknown quantum state.  Acquiring a
distinguishing measurement of \(|\psi\rangle\) requires either
declaring a measurement basis (with collapse) or accepting the
fidelity bound of an approximate cloner --- the \Fref{F12} apparatus
cost in either case
\citep{WoottersZurek1982NoCloning}.

\item[\emph{Goldwasser--Micali semantic security (cryptography).}] \textit{Analogy.}
For an IND-CPA secure probabilistic public-key encryption scheme,
no probabilistic polynomial-time adversary holding only the public
key and a challenge ciphertext can distinguish encryptions of two
adversarially chosen messages with non-negligible advantage.
Producing such a distinction requires apparatus extension --- the
secret key, a decryption oracle, or super-polynomial computation
breaking the underlying trapdoor (quadratic residuosity, factoring)
\citep{GoldwasserMicali1984ProbabilisticEncryption}.

\item[\emph{Myerson revenue equivalence (mechanism design).}] \textit{Analogy.}
In a symmetric independent-private-values auction, Myerson's envelope-formula characterisation shows that in every
incentive-compatible mechanism the allocation rule fixes each
bidder's interim expected payment up to a constant, the payment of
the lowest type, so mechanisms with the same allocation rule and
lowest-type payments raise the same expected revenue; any
incentive-compatible mechanism whose allocation responds to private
values leaves the bidders an information rent.  An allocation-only mechanism with no transfers collapses
to type-pooling
\citep{Myerson1981OptimalAuctionDesign}.
\end{description}
\endgroup

\subsection*{\Fref{F13a} Hiddenness Normal Form}
\begingroup\small
\begin{description}
\item[\emph{Pearl back-door criterion and \(do\)-calculus (causal
inference).}] \textit{Analogy.}
In a causal DAG with observed variables \(V\) and unobserved
confounders \(U\), an unblocked back-door path from \(X\) to \(Y\)
through a common ancestor in \(U\) is the canonical witness of a
hidden upstream determining state.  The observational/interventional
gap \(P(Y\mid X)\ne P(Y\mid \mathrm{do}(X))\) is then the named
predictive surplus inside a \(P(V)\)-fiber, and \(do\)-calculus
gives rules for deciding when causal effects remain nonparametrically identifiable, for instance via a back-door adjustment set or a front-door mediator (an instrumental variable alone generally does not identify the causal effect; additional identifying assumptions are needed, and these need not be parametric)
\citep{Pearl1995CausalDiagrams}.

\item[\emph{Le Verrier's prediction of Neptune (celestial
mechanics).}] \textit{Analogy.}
In 1845--1846, Le Verrier and (independently, from 1843) Adams analysed the residual
perturbations of Uranus's orbit after accounting for all
gravitational influences then catalogued.  The unexplained
predictive surplus in Uranus's observed position was interpreted as
the gravitational signature of an undiscovered upstream determining
state, and Le Verrier computed an ecliptic-longitude prediction
which Galle confirmed telescopically on 1846 September 23 within
\(1^\circ\) of Le Verrier's position: the lawful exposure of the
hidden planet Neptune
\citep{LeVerrier1846RecherchesUranus}.

\item[\emph{{\AA}str{\"o}m belief-MDP reduction for {POMDPs}
(reinforcement learning).}] \textit{Analogy.}
For a partially observable Markov decision process, two histories
sharing the current observation can yield different distributions
over future observation--reward sequences only when their belief
states (posteriors over the hidden Markov state) differ; distinct
belief states need not yield distinct distributions.
{\AA}str{\"o}m's 1965 reduction theorem realises the lawful exposure:
the POMDP is equivalent to a fully observed MDP on the
continuous-belief simplex \(\Delta(S)\), with the belief state as a sufficient upstream-quotient refinement,
not in general the minimal one, that resolves the
observation-fiber surplus
\citep{Astrom1965OptimalControlIncomplete}.

\item[\emph{Fisher's epistatic variance component (quantitative
genetics).}] \textit{Analogy.}
Under the Cockerham--Kempthorne orthogonal decomposition,
phenotypic variance partitions as
\(\operatorname{Var}(Y)=V_A+V_D+V_I+V_E\) with
\(V_I\) the epistatic interaction component.  A non-zero \(V_I\) is
the named witness of an upstream multi-locus determining state:
the phenotypic variability within a single-locus marginal fiber
that additive single-locus regression cannot explain corresponds
under the declared decomposition to hidden gene--gene interaction
structure
\citep{Fisher1918CorrelationRelatives}.
\end{description}
\endgroup

\subsection*{\Fref{F13b} CSL Formed-Closure Law}
In each item the determining readout \(d\) is the named hidden cause (assignable cause, implementation defect, orbiting body, resonance); the analogy asserts neither that \(d\) is determining nor that the cited source establishes it.
\begingroup\small
\begin{description}
\item[\emph{Shewhart statistical process control (industrial
quality engineering).}] \textit{Analogy.}
In a formed manufacturing closure with declared in-control mean
\(\mu\), standard deviation \(\sigma\), and \(3\sigma\) control
limits calibrated from a baseline period, an out-of-control signal
(point outside the control limits or a non-random run pattern) is
the named predictive surplus that triggers root-cause investigation
into an assignable cause --- tool wear, raw-material change,
operator effect, environmental drift --- under declared sampling
assumptions (a sampling rule, baseline independence, and identical
distribution).  The control-chart signal is a witness, not
a proof: false alarms exist under the in-control distribution at the
declared rate
\citep{Shewhart1931EconomicControl}.

\item[\emph{Model checking against a temporal-logic specification
(formal verification).}] \textit{Analogy.}
For a finite-state model \(M\) and a temporal-logic
specification \(\varphi\) (branching-time CTL in the cited algorithm,
linear-time LTL in SPIN), an explicit-state or symbolic
model-checker produces a counterexample trace \(\tau\) when
\(M\not\models\varphi\) for a linear-time or universal property.  The trace is the named predictive surplus
witnessing a hidden upstream defect in the implementation that the
declared specification was meant to forbid, with lawful exposure via
the automated checker (SPIN, NuSMV, TLC for TLA+ specifications)
under declared model-checking assumptions (a closed system, finite
state, a fair scheduler, and sound temporal-logic semantics)
\citep{ClarkeEmersonSistla1986AutomaticVerification}.

\item[\emph{Exoplanet detection by radial velocity and transit
(observational astronomy).}] \textit{Analogy.}
For a formed exoplanet-detection closure with declared instrument
(spectrograph + Doppler pipeline; or photometric transit pipeline)
and declared sensitivity (Doppler precision in m\,s\(^{-1}\);
photometric precision in ppm), a periodic radial-velocity
oscillation or photometric dimming inconsistent with stellar
activity is the named predictive surplus witnessing a hidden
orbiting body.  Mayor and Queloz realised this lawful exposure with
the 1995 detection of 51 Pegasi b, the first confirmed exoplanet
around a Sun-like star
\citep{MayorQueloz199551Pegasi}.

\item[\emph{{LHC} Higgs-boson discovery as resonance detection
(high-energy physics).}] \textit{Analogy.}
For a formed LHC search closure with declared integrated luminosity,
detector acceptance \(\times\) efficiency, and background-only
prediction (Standard-Model processes excluding the hypothesised
Higgs), an observed event yield exceeding the background-only model
by \(\ge 5\sigma\) in a signal region is the named predictive
surplus witnessing a hidden upstream resonance / particle predicted
by the broader theory.  The ATLAS 2012 observation of a 125 GeV
resonance in declared diphoton and four-lepton mass windows realised
the lawful exposure, with follow-up spin-parity and branching-ratio
measurements identifying the Standard-Model Higgs boson
\citep{ATLAS2012HiggsObservation}.
\end{description}
\endgroup

\subsection*{\Fref{F24} Layer Multiplicity}
\begingroup\small
\begin{description}
\item[\emph{{ANSI} {SQL} transaction isolation levels (database
concurrency).}] \textit{Analogy.}
For concurrent database transactions, the isolation-level quotient
classifies the visibility of in-flight operations.  A transaction's read or write role can split that quotient,
producing a dirty read, a non-repeatable read, or a phantom.  Each
transaction declares in advance one of the four ANSI isolation levels
(read-uncommitted, read-committed, repeatable-read, serializable),
and each level prices a different set of these role-split
obstructions as excluded or tolerated; Berenson et al.\ show that the
ANSI anomaly definitions are ambiguous and miss further anomalies
such as lost updates and write skew
\citep{BerensonBernsteinGrayMeltonONeilONeil1995SQLIsolation}.

\item[\emph{Grammatical case-marking (structural linguistics).}] \textit{Analogy.}
In a case-marking language the morphological word-form quotient
classifies nouns by inflection.  When a syntactic role (subject,
object, oblique, possessor) splits that morphological quotient, the
language's case system resolves the obstruction by assigning a
named case --- nominative, accusative, ergative, absolutive,
dative, genitive, locative, instrumental --- from Blake's
typological inventory.  Different languages instantiate different
subsets of the inventory, but each role-split is closed by a named
case-formation category
\citep{Blake2001Case}.

\item[\emph{{OSI} seven-layer networking reference model (computer
networking).}] \textit{Analogy.}
In a networking stack the protocol-layer quotient classifies
communication functions.  When a new communication role splits the
current layer quotient, the ISO/IEC 7498-1 model resolves the
obstruction by placing the function at one of seven named layers
(physical, data-link, network, transport, session, presentation,
application), each licensing specific function types and obstructing
others
\citep{ISOIEC7498OSI}.

\item[\emph{Polymorphism taxonomy (programming-language theory).}] \textit{Analogy.}
In a typed language the type quotient classifies expressions by
their monomorphic type.  When a function's required role is
generic across multiple types --- a role-split obstruction in the
monomorphic quotient --- the Cardelli--Wegner taxonomy resolves
the obstruction by assigning the function to a named polymorphism
category: parametric (universal quantification), ad-hoc (overloading /
coercion), or inclusion (subtyping).  Each category
licenses a distinct layer-formation mechanism for the generic role
\citep{CardelliWegner1985UnderstandingTypes}.
\end{description}
\endgroup

\subsection*{\Fref{F6} No-Needles Stability}
\begingroup\small
\begin{description}
\item[\emph{Quantum relative-entropy monotonicity.}] \textit{Analogy.}
For any completely positive trace-preserving (CPTP) quantum channel
\(T\) and density operators \(\rho,\sigma\), the relative entropy
satisfies \(S(T(\rho)\,\|\,T(\sigma))\le S(\rho\,\|\,\sigma)\).  The
decomposition is the Stinespring dilation
\(T(\rho)=\operatorname{Tr}_E[V\rho V^{*}]\); the monotonicity
combines joint convexity of relative entropy (Lindblad's lemma) with
isometric embedding and partial trace.  The preorder is scalar
\(\le\) on entropy functionals, not the Loewner block ordering
\citep{Lindblad1975CompletelyPositiveMaps}.

\item[\emph{Csiszár \(f\)-divergence data-processing inequality
(information theory).}] \textit{Analogy.}
For any convex \(f:(0,\infty)\to\mathbb R\) with \(f(1)=0\), the
\(f\)-divergence
\(D_f(P\,\|\,Q)=\mathbb E_Q\!\bigl[f(\mathrm dP/\mathrm dQ)\bigr]\) is
monotone under every Markov kernel \(K\):
\(D_f(K\!\cdot\!P\,\|\,K\!\cdot\!Q)\le D_f(P\,\|\,Q)\).  The
monotonicity comes from the joint convexity of
\((p,q)\mapsto qf(p/q)\) plus Jensen on the kernel mixture; a general
\(f\)-divergence has no chain rule.  The preorder is scalar \(\le\) on a
divergence value rather than the Loewner cone
\citep{LieseVajda2006Divergences}.

\item[\emph{Strassen monotone-coupling stochastic dominance
(probability theory).}] \textit{Analogy.}
By Strassen's theorem, \(X\le_{\mathrm{st}}Y\) iff there is a
coupling \(\pi\) of \(\mu_X\) and \(\mu_Y\) supported on
\(\{(x,y):x\le y\}\); applying a stochastically monotone Markov
kernel \(K\) to both coordinates preserves the order, so
\(K\!\cdot\!\mu_X\le_{\mathrm{st}}K\!\cdot\!\mu_Y\).  The
decomposition is the coupling existence (a measure-theoretic object
on a product space); the preorder is stochastic dominance defined by
tail-probability inequalities, with no positive-semidefinite cone
involved \citep{Strassen1965ExistenceProbabilityMeasures}.

\item[\emph{Lyapunov's direct method (dynamical systems).}] \textit{Analogy.}
For a dissipative dynamical system \(\dot x=f(x)\) with a
positive-definite Lyapunov functional \(V\) and a non-negative
dissipation
residual \(W\), the Lie-derivative inequality
\(\dot V=\nabla V\!\cdot\!f(x)\le-W(x)\le 0\) along trajectories
propagates \(V(x(t))\le V(x_0)\) for all \(t\ge 0\).  The
decomposition is the pair \((V,W)\) chosen against the vector field's
structure; transitivity composes through the flow semigroup
\(\varphi_t\), giving the \Fref{F6} endpoint along a continuous-time
flow rather than via a matrix Schur complement
\citep{Lyapunov1892StabilityMotion}.
\end{description}
\endgroup

\subsection*{\Fref{F8} Strict Extension}
\begingroup\small
\begin{description}
\item[\emph{Hematopoietic CMP \(\to\) MEP/GMP bifurcation (stem-cell
biology).}] \textit{Analogy.}
The common myeloid progenitor (CMP) population, prospectively
isolated by Lin\(^{-}\) Sca-1\(^{-}\) c-Kit\(^{+}\) CD34\(^{+}\)
Fc\(\gamma\)RII/III\(^{\mathrm{lo}}\) surface markers, splits into
megakaryocyte--erythroid (MEP) and granulocyte--macrophage (GMP)
progenitors distinguished by CD34 and Fc\(\gamma\)RII/III.  Both
daughters retain the Lin\(^{-}\) Sca-1\(^{-}\) c-Kit\(^{+}\) part of
the CMP gate and leave it on the CD34 and Fc\(\gamma\)RII/III axes
(MEP are CD34\(^{-}\) Fc\(\gamma\)RII/III\(^{\mathrm{lo}}\), GMP are
CD34\(^{+}\) Fc\(\gamma\)RII/III\(^{\mathrm{hi}}\)), and the witness
split is two
clonally related siblings of one CMP division that commit to MEP and
GMP respectively \citep{AkashiTraverMiyamotoWeissman2000CMPClonogenic}.

\item[\emph{Parisi replica-symmetry-breaking hierarchy (statistical
mechanics).}] \textit{Analogy.}
In Parisi's mean-field solution of the Sherrington--Kirkpatrick spin
glass, the one-step replica-symmetry-breaking ansatz partitions
overlaps between pure states into two values \(\{q_0,q_1\}\); full RSB
refines this to a continuous overlap function \(q(x)\) for
\(x\in[0,1]\).  Cutting the full-RSB ultrametric tree at one
overlap level groups pure states into clusters with the two-level
overlap pattern of the one-step ansatz; this cluster partition, not
the one-step ansatz solution itself, is the retention.  The
ultrametric witness split is a triple of pure states in one cluster,
all of whose pairwise overlaps lie above the cut, with
\(q_{\alpha\beta}\ne q_{\alpha\gamma}\) under full RSB
\citep{Parisi1983OrderParameterSpinGlasses}.

\item[\emph{Verner's Law refining Grimm's Law (historical
linguistics).}] \textit{Analogy.}
Grimm's Law sends PIE \({*}p,{*}t,{*}k,{*}k^w\) (outside clusters after \({*}s\) or, for \({*}t\), after another obstruent) to
Proto-Germanic \({*}f,{*}\theta,{*}h,{*}h^w\); on the reflexes of the PIE voiceless stops, Verner's Law refines
this output by voicing a non-initial fricative when the preceding PIE
syllable was unaccented.  The refinement holds only on that domain:
the voiced Verner reflexes merge with the reflexes of the PIE voiced
aspirates.  The
grammatischer-Wechsel alternation within Germanic strong-verb
paradigms is the canonical witness: forms that differed only in PIE accent position remain merged under
Grimm alone but split by voicing under Grimm + Verner, even though the conditioning accent is itself
lost in attested Germanic
\citep{Verner1877AusnahmeLautverschiebung}.
\end{description}
\endgroup

\subsection*{\Fref{F9} No Unstatused Residual}
\begingroup\small
\begin{description}
\item[\emph{Progress and preservation.}] \textit{Analogy.}
The Wright--Felleisen syntactic type-soundness theorem states that
every well-typed reachable configuration of a source term retains
its typing judgment (preservation) and is either a value or admits
a reduction step (progress).  The progress class
\(\{\text{value},\,\text{can-step}\}\) together with the retained
typing judgment is the \Fref{F9} typed status.  The orphan residual is a
\emph{stuck term}: a closed non-value with no reduction rule and no
type, forbidden by progress together with preservation
\citep{WrightFelleisen1994SyntacticTypeSoundness}.

\item[\emph{Genetic code via ribosomal translation (molecular
biology).}] \textit{Analogy.}
The ribosome decodes each mRNA codon into a unique amino acid (the
61 sense codons) or a termination signal (the 3 stop codons),
following the standard genetic code table; codon--anticodon pairing
at the A-site enforces single-valuedness, with context-dependent
recoding of UGA to selenocysteine and of UAG to pyrrolysine at
specific sites as the known exceptions.  Ribosomes stalled at the 3\('\) end of mRNAs lacking an in-frame stop codon constitute the translation-failure residual; tmRNA trans-translation in bacteria, and Pelota--Hbs1 splitting followed by ribosome-associated quality-control ubiquitination in eukaryotes, rescue the ribosome and mark the peptide for degradation\footnote{\url{https://doi.org/10.1126/science.271.5251.990}; \url{https://doi.org/10.1016/j.molcel.2014.07.006}.}
\citep{Crick1968GeneticCode}.

\item[\emph{Wick--Wightman--Wigner superselection (foundational
quantum physics).}] \textit{Analogy.}
The WWW superselection rules require that every physically
admissible pure state lie in a unique joint eigenspace of the
superselection observables --- univalence \((-1)^{2J}\) and electric
charge \(Q\).
Coherent superpositions across distinct sectors are excluded from
the physical state space because the relative phase is unobservable
(the superselection observables commute with every observable), so
the would-be superposition reduces to a sector-mixture rather than a
pure state
\citep{WickWightmanWigner1952IntrinsicParity}.

\item[\emph{Gale--Shapley deferred-acceptance matching
(mechanism design).}] \textit{Analogy.}
In a two-sided matching market with complete strict preferences and
equal market sides, the Gale--Shapley deferred-acceptance algorithm
terminates with a stable matching \(\mu\) that is a function: every agent is assigned exactly one partner
\citep{GaleShapley1962CollegeAdmissions}; later work showed that
truth-telling is a dominant strategy for the proposing side
\citep{DubinsFreedman1981Machiavelli,Roth1982EconomicsMatching}.  An
orphan residual, an agent with no partner \(\mu(x)\), cannot occur: that
every agent is matched plays the role of the \Fref{F9} coverage rule,
and that \(\mu\) is a function plays the role of agreement of statuses
at a common locus.
\end{description}
\endgroup

\subsection*{\Fref{F4} Local-Global Obstruction Normal Form}
\begingroup\small
\begin{description}
\item[\emph{Schema mapping in data integration (computer science).}] \textit{Analogy.}
In data integration, source instances over local schemas are
glued to a mediated global schema through declared
tuple-generating and equality-generating dependencies (tgds, egds).
Local-overlap compatibility is the egd-and-tgd satisfaction check
at the chase; egd failure in the chase signals that no solution exists; under weakly acyclic tgds the chase terminates and, when no egd fails, yields a universal solution, unique up to homomorphic equivalence (its core up to isomorphism), from which certain answers are computed \citep{Fagin2005DataExchange}.

\item[\emph{Genome assembly from overlapping reads (biology).}] \textit{Analogy.}
Sequencing reads are local sequence data on intervals of an
unknown genome; overlap-layout-consensus and de Bruijn graph
assemblers test pairwise overlaps for sequence agreement and
attempt to stitch compatible reads into a single contig.  The
four statuses appear as overlap conflicts, fragmented contigs
without a single traversal, repeat-induced chimeric-vs-unique
ambiguity, and unique reference assembly
\citep{PevznerTangWaterman2001Eulerian}.

\item[\emph{Pose-graph SLAM with loop closures (engineering).}] \textit{Analogy.}
A robot trajectory and map are estimated from local sensor
observations whose loop closures impose pairwise pose constraints.
The four statuses appear as chi-squared loop-closure rejection
(perceptual aliasing), nonzero cocycle drift around a cycle (no
global \(SE(3)\) embedding), gauge ambiguity (global rigid
transform, monocular scale), and converged g2o\,/\,GTSAM solution
with positive-definite information matrix
\citep{Cadena2016SLAMSurvey}.
\end{description}
\endgroup

\subsection*{\Fref{F14} Duality Fixity / Self-Dual Trace Confinement}
\begingroup\small
\begin{description}
\item[\emph{Bender--Boettcher {PT}-symmetric Hamiltonians
(foundational physics).}] \textit{Analogy.}
For a spin-\(0\) quantum system, the composite \(\mathcal{PT}\)
operator is antilinear with \((\mathcal{PT})^2=\mathrm{id}\), an
involution on the Hilbert space.  Bender and Boettcher conjectured, on numerical and WKB evidence (reality for \(\varepsilon\ge0\) was later proved by Dorey, Dunning and Tateo), that
\(\mathcal{PT}\)-symmetric non-Hermitian Hamiltonians (such as
\(H=p^2+x^2(ix)^{\varepsilon}\), including \(H=p^2+ix^3\) at
\(\varepsilon=1\)) have entirely real spectra precisely when the
\(\mathcal{PT}\) symmetry is unbroken: each physical eigenstate is
a \(\mathcal{PT}\)-fixed eigenstate; a nonzero imaginary part of an eigenvalue,
the \(\mathcal{PT}\)-skew spectral component, occurs only when the symmetry is
broken. The analogy is with confinement to a fixed locus; no domination ladder
occurs in the cited work
\citep{BenderBoettcher1998PTSymmetric}.

\item[\emph{Real signals as the Hermitian-symmetric fixed locus (signal
processing).}] \textit{Analogy.}
For a finite-energy signal \(x[n]\in\mathbb C\) over \(\mathbb Z\), the involution \(J(X)(\omega)=\overline{X(-\omega)}\) on its discrete-time Fourier transform has as fixed locus the Hermitian-symmetric spectra, those of real signals.  The
\(J\)-anti-invariant spectral component is the \Fref{F14} anti-invariant
readout, and vanishing of its integrated power forces the spectrum
into that fixed locus, that is, \(x\) real (by Parseval its power is the energy of \(\operatorname{Im}x\))
\citep{OppenheimSchafer2010DiscreteTimeSP}.

\item[\emph{Galois conjugation and the fixed-field theorem
(algebraic number theory).}] \textit{Analogy.}
For a Galois extension \(L/K\) with cyclic order-2 Galois group
\(\langle\sigma\rangle\), the fundamental theorem of Galois theory
identifies the base field \(K\) with the fixed locus \(L^{\sigma}\)
of the involution \(\sigma\).  The anti-invariant component
\(z-\sigma(z)\) is the \Fref{F14} anti-invariant readout, and \(z-\sigma(z)=0\) forces
\(z\in K\) --- confinement to the fixed field
\citep{Galois1846Memoire}.

\item[\emph{Smith theory: \(\mathbb Z/2\)-cohomology and
fixed-point sets (algebraic topology).}] \textit{Analogy.}
P.~A. Smith's theorem relates the
\(\mathbb Z/2\)-cohomology of a topological space \(X\) under a
continuous involution \(J:X\to X\) to its fixed-point set
\(X^J\).  For a finite \(\mathbb Z/2\)-CW complex \(X\),
\(\dim H^*(X^J;\mathbb Z/2)\le\dim H^*(X;\mathbb Z/2)\) (the Smith--Floyd inequality), with
equality when \(X\) is totally non-homologous to zero.  With \(\mathbb Z/2\) coefficients the sign change of \(J\) is
invisible, so the \Fref{F14} anti-invariant readout has only a
counterpart here: Borel localisation, under which the equivariant
cohomology \(H^*_{\mathbb Z/2}(X;\mathbb Z/2)\), after inverting the
generator of \(H^1(B\mathbb Z/2;\mathbb Z/2)\), is carried by the
fixed-point set
\citep{Smith1938TransformationsFinitePeriod}.
\end{description}
\endgroup

\subsection*{\Fref{F15a} Selection Obstruction / Repair Normal Form}
\begingroup\small
\begin{description}
\item[\emph{Hairy ball theorem / Poincar{\'e}--Hopf (differential
topology).}] \textit{Analogy.}
On the 2-sphere \(S^2\), no continuous nonvanishing tangent vector
field exists.  The Poincar{\'e}--Hopf theorem identifies the
obstruction with the Euler characteristic
\(\chi(S^2)=2\neq 0\), and it rules out every nonvanishing field,
equivariant or not.  \(\mathrm{SO}(3)\)-equivariance is more
restrictive still: the rotations fixing a point fix no nonzero
tangent vector there, so the zero field is the only equivariant
section of \(TS^2\).  A declared breaker --- a preferred polar axis
or a stochastic perturbation --- selects a non-equivariant field,
which still vanishes somewhere; a nonvanishing direction field exists
only after points are removed, such as the poles of the declared
axis
\citep{Poincare1885CourbesEquationsDifferentielles}.

\item[\emph{Schelling focal points in coordination games (game
theory).}] \textit{Analogy.}
In a pure-coordination game with multiple symmetric equilibria, no
symmetric pure strategy can pin down a unique equilibrium.  Schelling
identified \emph{focal points} --- salient landmarks, precedents,
or conventions distinguishing one equilibrium from its symmetric
partners --- as the \Fref{F15a} declared breaker; mixed-strategy
randomisation provides an equivariant but non-selective
alternative
\citep{Schelling1960StrategyConflict}.
\end{description}
\endgroup

\subsection*{\Fref{F15b} Spontaneous Symmetry Breaking / Branch-Selection Formation}
\begingroup\small
\begin{description}
\item[\emph{Rayleigh--B{\'e}nard convection (fluid dynamics).}] \textit{Analogy.}
A horizontal fluid layer heated from below transitions, at a
critical Rayleigh number (\(\mathrm{Ra}_c\approx 1708\) for the
standard rigid no-slip boundary case and lower for stress-free
boundaries), from a translation- and rotation-invariant
conduction state to a patterned state of convection rolls or
hexagonal cells.  Each experimental realisation fixes a particular
roll axis or hexagonal sublattice, selecting a branch in the
continuous symmetry orbit
\citep{Chandrasekhar1961HydrodynamicStability}.

\item[\emph{Soai chiral autocatalysis (chemistry).}] \textit{Analogy.}
In the Soai reaction the product catalyses its own formation with
the same chirality; starting from a slightly enantioenriched mixture, autocatalysis amplifies the initial bias to nearly enantiopure product; started without chiral bias, either enantiomer results stochastically from run to run (reported by Soai and coworkers in 2003\footnote{\url{https://doi.org/10.1016/S0957-4166(02)00791-7}.}).  The underlying kinetics are chirally symmetric, but the
autocatalytic formation operator picks a single enantiomer per
realisation
\citep{SoaiShibataMoriokaChoji1995AsymmetricAutocatalysis}.

\item[\emph{Biological homochirality of life (origin-of-life
biology).}] \textit{Analogy.}
Terrestrial biology employs L-amino acids and D-sugars almost
exclusively, despite a mirror-symmetric prebiotic chemistry.
Candidate \Fref{F15b} formation-mechanism classes include
Frank-style autocatalytic amplification, photochemical bias from
circularly polarised light, and evolutionary locking after an
early enantiomeric bias; each names a declared branch-selection
source by which the biosphere settled into one chirality branch
\citep{Bonner1991OriginAmplificationBiomolecularChirality}.

\item[\emph{Crystallisation from supercooled liquid (condensed
matter).}] \textit{Analogy.}
A homogeneous, isotropic liquid below its freezing point is
thermodynamically metastable; nucleation of a crystalline phase
breaks continuous translations and \(\mathrm{SO}(3)\) rotations to
a discrete crystallographic subgroup, with each nucleation event
selecting a particular lattice orientation and translation phase
\citep{Turnbull1950FormationCrystalNuclei}.
\end{description}
\endgroup

\subsection*{\Fref{F40} Anomaly / Symmetry Obstruction}
\begingroup\small
\begin{description}
\item[\emph{Gimbal lock in Euler-angle parametrisations
(3D graphics).}] \textit{Analogy.}
The Lie group \(\mathrm{SO}(3)\) is a smooth 3-manifold, but its
Euler-angle parametrisation has coordinate singularities --- gimbal
lock --- where two rotation axes align.  At those configurations
DomainDescends fails: the parametrisation domain breaks down and
the orientation readout loses a degree of freedom.  Quaternion or
rotation-matrix representations are the declared apparatus repair
\citep{Kuipers1999QuaternionsRotationSequences}.

\item[\emph{Discretized translation symmetry.}] \textit{Analogy.}
A continuum PDE such as Navier--Stokes has continuous translational
symmetry, and its momentum equation is in flux form, so total momentum is conserved when net boundary momentum flux and net external force vanish (for the inviscid Euler equations this is the Noether charge of translation invariance).  A finite-difference discretisation on a fixed spatial
lattice respects only a discrete translation subgroup: both
SymmetryDescends and DomainDescends fail for the continuous group.  A
conservative flux-form discretisation still conserves discrete total
momentum exactly; a non-conservative one does not preserve the
downstream momentum readout.  Flux-form spatial discretisations,
paired in time with symplectic or energy-momentum-preserving
variational integrators, are the apparatus repair
\citep{HairerLubichWanner2006GeometricNumerical}.

\item[\emph{Triangular arbitrage in currency exchange (finance).}] \textit{Analogy.}
In a frictionless ideal market, exchange rates would compose
multiplicatively to a triangular product equal to 1
(e.g.\ \(\mathrm{USD/EUR}\cdot\mathrm{EUR/JPY}\cdot
\mathrm{JPY/USD}=1\)).  Quotes that update asynchronously across currency pairs break
GroupLawPreserved on the round-trip composition, and a product
different from unity offers a triangular-arbitrage opportunity in one
direction around the triangle; bid--ask spreads and transaction fees
lower the round-trip product in both directions and so absorb small
inconsistencies.  Competitive-arbitrage pressure, which restores
consistent quotes, and explicit modelling of friction are the
declared apparatus repairs
\citep{Aliber1973InterestRateParity}.

\item[\emph{Galilean vs special-relativistic velocity addition
(foundational physics).}] \textit{Analogy.}
The would-be Galilean velocity-addition group \(v_{AB}=v_{AC}+v_{CB}\)
fails to be closed on the physical velocity domain \(|v|<c\):
DomainDescends fails when Galilean composition of sublight
velocities can exit the declared domain, and the velocity readout
fails to be preserved.  Einstein's relativistic addition for collinear velocities,
\(v_{AB}=(v_{AC}+v_{CB})/(1+v_{AC}\cdot v_{CB}/c^2)\) --- equivalently
rapidity additivity --- is the declared apparatus repair
\citep{Einstein1905Elektrodynamik}.
\end{description}
\endgroup

\subsection*{\Fref{F41} Duality Equivalence}
\begingroup\small
\begin{description}
\item[\emph{Stone duality (logic and topology).}] \textit{Analogy.}
Stone's representation theorem gives a contravariant equivalence
between Boolean algebras and Stone spaces (compact totally
disconnected Hausdorff spaces): \(\mathrm{Spec}\) sends a Boolean
algebra to its ultrafilter space, \(\mathrm{Clop}\) sends a Stone
space to its Boolean algebra of clopen subsets, and
\(A\cong\mathrm{Clop}(\mathrm{Spec}(A))\),
\(X\cong\mathrm{Spec}(\mathrm{Clop}(X))\).  Role readouts (Boolean
operations \(\leftrightarrow\) topological operations on clopens),
routes (homomorphisms \(\leftrightarrow\) continuous maps), and
status (the Boolean-algebra / Stone-space structures) are
preserved on both sides
\citep{Stone1936RepresentationsBooleanAlgebras}.

\item[\emph{Gelfand duality (operator algebras and topology).}] \textit{Analogy.}
The Gelfand--Naimark theorem gives a contravariant equivalence
between commutative unital \(C^*\)-algebras and compact Hausdorff
spaces: \(\mathrm{Spec}\) returns the space of characters and
\(C(\cdot,\mathbb C)\) returns the algebra of continuous functions,
with the Gelfand transform an isometric \(*\)-isomorphism
\(A\cong C(\mathrm{Spec}(A))\) and homeomorphism
\(X\cong\mathrm{Spec}(C(X))\).  Role readouts (\(*\)-algebra
operations \(\leftrightarrow\) continuous-function operations),
routes (\(*\)-homomorphisms \(\leftrightarrow\) continuous maps),
and status (commutative-\(C^*\) versus compact-Hausdorff) are
preserved
\citep{GelfandNaimark1943NormedRings}.

\item[\emph{Curry--Howard correspondence (logic and type theory).}] \textit{Analogy.}
For the implicational fragment of intuitionistic natural deduction
and the simply typed \(\lambda\)-calculus (with product, sum and
empty types added for full propositional logic), Curry--Howard gives an exact syntactic
equivalence between propositions and types, proofs and
\(\lambda\)-terms: introduction/elimination rules correspond to
abstraction/application, and detour
elimination (proof normalisation) corresponds to
\(\beta\)-normalisation.  Roles (formula \(\leftrightarrow\) type,
proof \(\leftrightarrow\) program), routes (proof rules
\(\leftrightarrow\) typing rules), and status (provability
\(\leftrightarrow\) inhabitation, normalisation) are preserved by
mutually inverse translations between the two formal systems
\citep{Howard1980FormulaeAsTypes}.

\item[\emph{{AdS}/{CFT} holographic correspondence (string
theory; conjectural).}] \textit{Analogy.}
Maldacena's AdS/CFT correspondence conjectures that quantum gravity
on a \((d+1)\)-dimensional anti-de Sitter spacetime is dual to a
conformal field theory on its \(d\)-dimensional asymptotic
boundary, with the canonical example type IIB string theory on
\(\mathrm{AdS}_5\times S^5\) and \(\mathcal N=4\) super-Yang--Mills
on the boundary.  Within the accepted holographic dictionary, bulk
gravitational operators map to boundary CFT operators, partition
functions match, and AdS isometries map to CFT conformal
symmetries --- the proposed \Fref{F41} role/route/status preservation
across the duality
\citep{Maldacena1997LargeN}.
\end{description}
\endgroup

\subsection*{\Fref{F17} Objectivity as Common Quotient}
\begingroup\small
\begin{description}
\item[\emph{State-machine replication via Paxos (distributed
systems).}] \textit{Analogy.}
In a Paxos-style replicated state machine, each replica observes
its own log of accepted operations while the consensus protocol
commits a single sequence that all non-faulty replicas converge
on.  The per-replica log views are \Fref{F17} access quotients, the
committed sequence is the public quotient, and the safety property
of consensus is the common-across-accesses condition; split-brain
configurations are the \Fref{F17} access-objectivity defect
\citep{Lamport1998PartTimeParliament}.

\item[\emph{The {WGS-84} geodetic reference frame (geodesy).}] \textit{Analogy.}
National and regional geodetic datums (NAD-83, ETRS-89, JGD2011)
each provide a local access quotient on geodetic position, and
each maps to WGS-84 via published transformations.  WGS-84 is the
\Fref{F17} public quotient through which the local datums factor, and
agreement under WGS-84 reduction is the common-across-accesses
property
\citep{NGA2014WGS84}.

\item[\emph{The {SI} as public measurement quotient (metrology).}] \textit{Analogy.}
National metrology institutes (NIST, NPL, PTB, etc.) maintain
their own primary standards realising the seven SI base units
through the seven defining constants
\((c,h,e,k,N_A,K_{\mathrm{cd}},\Delta\nu_{\mathrm{Cs}})\).  The SI
is the public quotient through which all national standards
factor, and SI-traceable measurement agreement across laboratories
under the CIPM Mutual Recognition Arrangement is the \Fref{F17}
common-across-accesses property; unresolved Key-Comparison
discrepancies are the access-objectivity defect
\citep{BIPM2019SIBrochure}.

\item[\emph{Thermodynamic ensemble equivalence (statistical
mechanics).}] \textit{Analogy.}
In the thermodynamic limit and for short-range systems away from
phase transitions, the microcanonical, canonical, and grand-canonical
ensembles yield equivalent equations of state for macroscopic
intensive quantities.  Each ensemble is an access quotient on
phase space, and the thermodynamic equation of state is the \Fref{F17}
public quotient through which all three factor.  Ensemble
inequivalence at long-range systems / phase-transition
pathologies is the access-objectivity defect
\citep{Ruelle1969StatisticalMechanics}.
\end{description}
\endgroup

\subsection*{\Fref{F18} Lawhood Descent}
\begingroup\small
\begin{description}
\item[\emph{Stone--{\v C}ech compactification and continuous
extensions (topology).}] \textit{Analogy.}
For a Tychonoff space \(X\), the Stone--{\v C}ech
compactification \(\beta X\) is characterised by the universal
property that every continuous bounded real-valued function on
\(X\) extends uniquely to a continuous function on \(\beta X\).
A predicate declared on a dense subset is lawful on the closure
iff it admits a continuous extension --- the \Fref{F18} descent
criterion for topological laws
\citep{Cech1937BicompactSpaces}.
\end{description}
\endgroup

\subsection*{\Fref{F19} Object Persistence}
\begingroup\small
\begin{description}
\item[\emph{Comoving galaxy identity in FLRW cosmology
(cosmology).}] \textit{Analogy.}
In a Friedmann--Lema{\^\i}tre--Robertson--Walker universe, galaxies
moving with the Hubble flow keep their comoving coordinates as the
scale factor evolves.  Comoving-coordinate identity persists under
the cosmic-time transport; peculiar velocities and gravitational
mergers are the persistence obstruction
\citep{Friedmann1922KrummungRaumes}.

\item[\emph{Round-trip identity under {JSON} serialisation
(computer science).}] \textit{Analogy.}
For a data value \(D\), the round-trip
\(D\equiv\mathrm{deserialise}(\mathrm{serialise}(D))\) is the \Fref{F19}
persistence condition for the declared canonical-form identity
quotient.  Locale-dependent number formatting, time-zone-aware
datetimes, hash-randomised dictionary ordering, and floating-point
precision loss are common persistence obstructions
\citep{Crockford2006JSON}.

\item[\emph{Corporate personhood across statutory merger
(corporate law).}] \textit{Analogy.}
Under the Model Business Corporation Act \S\,11.07, a statutory
merger transfers the corporate personhood of constituent
corporations to a single surviving (or new) corporation, which
inherits the rights, properties, and obligations of the
constituents.  The legal-identity transport descends through the
corporate-personhood quotient under declared statutory procedures;
defective merger filings or undocumented dissolution without
successor are the persistence obstruction
\citep{MBCA2016Section1107}.
\end{description}
\endgroup

\subsection*{\Fref{F20} Fact / Record Formation}
\begingroup\small
\begin{description}
\item[\emph{Blockchain ledger and confirmation depth (distributed
computer science).}] \textit{Analogy.}
In a blockchain, each transaction is recorded by cryptographic
signature, included in a block, and chained to prior blocks via
hash pointers; ``stable fact'' is properly indexed by
confirmation depth.  Signature-verification failure is the
event-recording obstruction, hash tampering or chain
reorganisations and forks are the record-stability obstruction,
and a double-spend constitutes the fact-value obstruction
\citep{Nakamoto2008Bitcoin}.

\item[\emph{Forensic chain of custody (criminal procedure).}] \textit{Analogy.}
Physical evidence is admissible only when its chain of custody
is fully documented from collection through analysis: each
transfer between custodians is the recorded event, an unbroken
custodial log is the stable record, and laboratory results
matching the seized item are the consistent fact-value.  Gaps,
mishandling, or tampering are \Fref{F20} stable-fact failures
\citep{Rice2008ElectronicEvidence}.

\item[\emph{Pacioli double-entry bookkeeping (accounting).}] \textit{Analogy.}
Pacioli's double-entry method requires every transaction to be
posted as two equal opposing entries, preserving the accounting
identity \(\text{Assets}=\text{Liabilities}+\text{Equity}\).  An
unrecorded transaction is the event-recording obstruction, a
posting error is the record-stability obstruction, and a
trial-balance discrepancy is the fact-value obstruction
\citep{Pacioli1494Summa}.

\item[\emph{Paleontological taphonomy and the fossil record
(natural science).}] \textit{Analogy.}
Fossil-record formation is the natural-history analogue of \Fref{F20}:
preserved burial events, geological persistence, and consistent
inferential analyses yield robust paleontological facts.  Pre-burial
scavenging and decay are the event-recording obstruction;
metamorphism, erosion, and exhumation damage are the
record-stability obstruction; conflicting species-identification
or dating analyses are the fact-value obstruction.  The fossil
record is therefore selective rather than complete, with stable
facts indexed by preservation-and-inferential quality
\citep{KidwellHolland2002QualityFossilRecord}.
\end{description}
\endgroup

\subsection*{\Fref{F21} Reflexive Nonclosure}
\begingroup\small
\begin{description}
\item[\emph{Sarbanes-Oxley independent-audit mandate
(corporate-finance regulation).}] \textit{Analogy.}
Since the Securities Act of 1933 and the Securities Exchange Act of
1934, public-company financial statements have been audited by an
independent public accountant rather than by the issuer's own
accounting staff; the Sarbanes-Oxley Act of 2002 tightened
auditor-independence rules and added internal-control reporting.
Management's internal-controls assessment under \S404(a) is the
scoped lower audit; the independent auditor's attestation under
\S404(b), required of accelerated filers, is the meta-layer
certificate, and the auditor is itself overseen by the Public Company
Accounting Oversight Board that the Act created, producing a tower of
meta-audits.  Complete
self-attestation by the issuer is outside the scope of
the regime \citep{SarbanesOxley2002Act}.

\item[\emph{X.509 PKI external root-of-trust
(cryptographic protocol design).}] \textit{Analogy.}
In the X.509 public-key infrastructure, an endpoint
certificate cannot self-certify its trustworthiness to a
relying party: validation requires a chain of signatures
terminating at a root certificate authority in the
client's externally maintained trust store, which acts
as the meta-layer.  Checks on the leaf alone (syntax, validity period, key
usage) are the scoped lower audit; the leaf carries its issuer's
signature, not its own, and self-signed leaves are rejected by default by
mainstream TLS clients.  Complete in-certificate
self-trust is outside the scope of the protocol
\citep{Cooper2008RFC5280}.

\item[\emph{Rice's theorem on undecidability of semantic
properties (computability theory).}] \textit{Analogy.}
Rice's theorem rules out any same-layer universal
semantic decider for non-trivial behavioural properties
of partial recursive functions; specialised to
self-application, no program can act as a complete
same-layer
semantic self-certifier for a non-trivial behavioural
property.  Bounded-resource testing on a finite input
set within a step bound is a scoped lower audit, accepted
but insufficient for the full semantic claim.  Legitimate
reflexive certification of a particular instance must
invoke a strictly stronger meta-formalism (proof-assistant
formalisation against a specification, abstract
interpretation sound for the target property, or
model-checking of a decidable abstraction), which
certifies particular claims rather than serving as a
same-layer universal decider \citep{Rice1953Classes}.

\item[\emph{M\"unchhausen trilemma in epistemology
(analytic epistemology).}] \textit{Analogy.}
Hans Albert's M\"unchhausen trilemma argues that any
attempted self-grounding of a knowledge claim faces
three exhaustive options, none of which constitutes
complete same-layer self-certification: infinite
regress, logical circularity, or arbitrary termination
in unargued axioms.  A locally consistent dialectical
chain of reasons is the scoped lower audit; a legitimate
epistemic certificate must invoke a meta-layer
commitment (foundationalist intuition, coherentist
network, externalist reliability, or transcendental
argument), each of which exits the same-layer chain of
reasons \citep{Albert1985TreatiseCriticalReason}.
\end{description}
\endgroup

\subsection*{\Fref{F22} Interface Mediation}
\begingroup\small
\begin{description}
\item[\emph{Pharmacophore / structure-activity relationship
(medicinal chemistry).}] \textit{Analogy.}
Within the declared pharmacophore model of a target
receptor, the source quotient collapses candidate ligands
to the spatial / electronic arrangement of features
required for binding (hydrogen-bond donors and acceptors,
hydrophobic regions, charged groups); the boundary is the
receptor binding-site geometry; the update is binding
and downstream signalling; the target quotient is the
pharmacological response class.  Two molecules sharing
the same pharmacophore are predicted, within the model,
to produce the same response class at the same receptor;
deviations attributable to off-pharmacophore substituents
(allosteric effects, distinct ADMET profiles) are
diagnosed as boundary-leakage corrections outside the
declared interface \citep{Wermuth1998IUPACGlossary}.

\item[\emph{Vienna Convention on Diplomatic Relations
(international law).}] \textit{Analogy.}
The 1961 Vienna Convention formalises the diplomatic
interface between sovereign states.  The carrier is the
sending state's internal political and administrative
state; the source quotient collapses internal states
producing the same officially endorsed position; the
boundary is the recognised diplomatic channel (note
verbale, official meeting, protected correspondence)
delivered via accredited diplomatic agents; the update
is the receiving state's policy response through its
foreign ministry; the target quotient is the receiving
state's diplomatic posture.  Article 41 explicitly
forbids diplomats from interfering in the receiving
state's internal affairs, codifying the prohibition on
source-internal leakage past the diplomatic boundary
\citep{ViennaConventionDiplomatic1961}.

\item[\emph{Thermodynamic state functions and boundary
protocol (classical thermodynamics).}] \textit{Analogy.}
For a thermodynamic system following a declared
boundary protocol (e.g., quasi-static, adiabatic, or
isothermal exchange under specified constraints on
volume, pressure, and particle exchange), the carrier
is the microstate space; the source quotient is the
macroscopic state (energy, volume, particle number);
the boundary is the protocol-fixed heat and work
exchange together with its constraint structure; the
update is microscopic dynamical evolution; the target
quotient reads the post-process macrostate.  Within the
declared protocol, the macrostate change depends only
on the initial macrostate and the protocol-fixed
boundary data, not on the microstate trajectory; this is macroscopic determinism,
the premise of a thermodynamic description, and the state functions
\(U\) and \(S\) are the readouts whose changes depend only on the
initial and final equilibrium macrostates, not on the process path \citep{Callen1985Thermodynamics}.
\end{description}
\endgroup

\subsection*{\Fref{F25} Causality as Stable Intervention Transport}
\begingroup\small
\begin{description}
\item[\emph{Mendelian randomisation
(instrumental-variable epidemiology).}] \textit{Analogy.}
A germline genetic variant associated with a modifiable
exposure is used as a natural instrumental intervention
to estimate the causal effect of the exposure on a
disease outcome.  The carrier is the study population;
the context quotient collapses subjects sharing the
same confounder profile; the intervention space ranges
over genotypes at the chosen instrument locus; the
abstraction collapses genotypes to their associated
exposure level; the transport is the lifelong exposure
process; the readout is the disease outcome.
Admissibility rests on the conjunction of relevance,
independence (incorporating population-stratification
control), and the exclusion restriction, together with monotonicity (or effect homogeneity); the descended
residue is then the local-average exposure-outcome causal
effect \citep{DaveySmith2003Mendelian}.

\item[\emph{CRISPR-Cas9 directed gene knockout
(molecular biology).}] \textit{Analogy.}
An engineered guide-RNA-directed nuclease introduces a
targeted double-strand break at a specified genomic
locus, producing loss-of-function alleles in an
otherwise isogenic background.  The carrier is the
population of treated cells or organisms; the context
quotient collapses cells to their genetic background and
cell type; the intervention space ranges over candidate
guide RNAs and delivery conditions; the abstraction
collapses guide RNAs to the target locus they edit; the
transport is the editing event plus downstream
biological development; the target readout is the
phenotypic measurement.  Within the declared
editing-efficiency and off-target-control regime, the
phenotypic readout descends through
(background-and-cell-type, target-locus)
\citep{Jinek2012CRISPR}.

\item[\emph{Particle-scattering cross-section
measurement (high-energy physics).}] \textit{Analogy.}
A scattering experiment fires a beam of particles of
defined type and energy at a target of defined
composition; the carrier is the ensemble of individual
beam-target interaction events; the context quotient
collapses events to the target-and-beam apparatus
configuration; the intervention space parameterises
beam-energy and target choices; the abstraction
collapses operational details to the beam-target class;
the transport is the scattering process; the target
readout of a single event is the realised scattering
angle and energy.  Individual events are stochastic; the
\Fref{F25} reading is its probabilistic form, with the
differential cross-section \(d\sigma/d\Omega\) as the
stable kernel residue over repeated events.
The alpha-particle scattering off thin foils measured by Geiger and
Marsden from 1909 to 1913, with Rutherford's 1911 analysis, is the
founding instance
\citep{GeigerMarsden1913Deflexion,Rutherford1911Scattering}.

\item[\emph{Regression-discontinuity natural experiment
(empirical econometrics).}] \textit{Analogy.}
A regression-discontinuity design uses a known
deterministic cutoff in an assignment variable (test
score for scholarship eligibility, age for retirement
benefit, vote share for incumbency) to identify the
local average treatment effect at the threshold.  The
carrier is the population of units near the threshold;
the context quotient collapses units sharing the same
local cutoff neighbourhood of the running variable; the
intervention space is the treatment-assignment
indicator; the abstraction collapses delivery-channel
details to the just-above / just-below assignment; the
transport is the policy delivery process; the target
readout is the outcome of interest.  Admissibility
holds in the limit at the threshold; residue stability
is established by continuity of potential outcomes and
no-manipulation of the running variable around the
cutoff \citep{ThistlethwaiteCampbell1960Regression}.
\end{description}
\endgroup

\subsection*{\Fref{F26} Moduli / Contingency}
\begingroup\small
\begin{description}
\item[\emph{Instruction set architecture versus
micro-architecture (computer architecture).}] \textit{Analogy.}
The carrier is the space of computer-system
implementations; the lawfulness predicate requires
correct execution of the declared instruction set
architecture (ISA); the equivalence is observational
equivalence at the ISA boundary (same program, same
architectural-state update); the moduli quotient sends
an implementation to its ISA-conformant behaviour
class.  The ISA contract (instruction encodings,
register semantics, addressing modes, observable side
effects) is a necessary readout shared by all
conforming implementations.  The micro-architecture
(pipeline depth, cache hierarchy, branch-prediction
strategy, clock frequency) is presentation-dependent: it
varies among implementations of one behaviour class.  The
System/360 family, which fixed an ISA while shipping
many distinct micro-architectures, is the founding
instance \citep{AmdahlBlaauwBrooks1964System360}.

\item[\emph{Evolutionary contingency in macroevolution
(paleobiology).}] \textit{Analogy.}
The carrier is the space of possible biotic
trajectories on Earth; the lawfulness predicate
requires consistency with natural selection,
developmental constraints, and the physical
environment; the equivalence identifies trajectories
with identical sequences of macroevolutionary
bottleneck outcomes; the moduli quotient sends a
trajectory to its bottleneck-survival-pattern class.
Selection-driven adaptation, the existence of
developmental constraints, and thermodynamically
consistent matter/energy flow are necessary readouts.
The specific assemblage of surviving phyla on Earth
(e.g., the contingent survival of chordates after the
Cambrian extinction events) is a contingent readout
that varies across moduli classes; ``replaying the
tape of life'' would traverse a different moduli class
\citep{Gould1989WonderfulLife}.
\end{description}
\endgroup

\subsection*{\Fref{F28} Composability}
\begingroup\small
\begin{description}
\item[\emph{Retrosynthetic analysis and synthon-based
total synthesis (organic chemistry).}] \textit{Analogy.}
The carriers are organic compounds; the source quotients
collapse compounds to their synthon class (connectivity
and functional groups); the interface readouts are
reactive sites available for bond formation; the
composition is a declared coupling reaction within a
specified reaction family of declared functional-group
tolerance.  The residual records yield, regio- and
stereo-byproduct distribution, and
reaction-mass-efficiency residual.  Composite-outcome
descent holds
when the same target connectivity is obtained from
distinct yet retrosynthetically-equivalent precursors;
cross-residual descent holds when the residual
byproduct distribution is invariant under
representative-level choices within the same reaction
family \citep{Corey1989LogicSynthesis}.

\item[\emph{Intermodal maritime containerisation
(industrial logistics).}] \textit{Analogy.}
The ISO 668 standard fixes the external dimensions,
corner-fitting geometry, and rated load of intermodal
freight containers.  The carriers are containers loaded
by different shippers; the source quotients collapse
them to their ISO 668-conformant external-interface
class; the interface readouts are the ISO corner
fittings and the rated gross mass; the composition is
stacking on a vessel or chaining across
truck-train-vessel handoffs; the composite quotient
classifies the consolidated load by its
external-interface footprint.  The residual captures
internal load distribution, weight class, and
damage / customs / handling residuals.  Admissibility
restricts to rated compatible cargo (weight under
ratings, stacking-strength conformance, hazmat
compatibility) \citep{ISO668_2020}.

\item[\emph{{OSI} 7-layer protocol stack composition
(network engineering).}] \textit{Analogy.}
The OSI Basic Reference Model partitions network
communication into seven layers, each defined by
service primitives at Service Access Points (SAPs).
The carriers are layer-\(N\) entity implementations;
the source quotients collapse implementations to their
SAP-service-primitive contract; the interface readouts
record the service-primitive sequence delivered at the
SAP; the composition stacks a layer-\((N{+}1)\) entity
on top of a layer-\(N\) service by invoking layer-\(N\)
service primitives; the composite quotient classifies
the composite by its end-to-end service.  The residual
records PDU header overhead, fragmentation, and
retransmission residual.  Composite-outcome descent
holds because the layer-\((N{+}1)\) service depends
only on layer-\(N\) service primitives, not on the
specific implementation of layer \(N\)
\citep{ISOIEC7498OSI}.

\item[\emph{Golden Gate modular {DNA} cloning
(synthetic biology).}] \textit{Analogy.}
Golden Gate cloning uses Type IIS restriction enzymes
that cut DNA outside their recognition site, leaving
scarless four-base overhangs designed arbitrarily.
The carriers are donor plasmids carrying DNA parts
flanked by Type IIS recognition sequences; the source
quotients collapse plasmids to their part-class plus
overhang specification; the interface readouts are the
four-base overhang fingerprints; the composition is the
one-pot enzymatic digest-and-ligate reaction; the
composite quotient classifies the assembled multi-part
construct.  The residual records misligation products,
unligated linearised backbone, and off-target
re-religation events.  Composite-outcome descent holds
because overhang complementarity determines assembly
order independent of donor-plasmid backbone
\citep{Engler2008GoldenGate}.
\end{description}
\endgroup

\subsection*{\Fref{F30} Explanation as Compression with Descent}
\begingroup\small
\begin{description}
\item[\emph{Mendeleev periodic table explaining
chemical regularities (chemistry).}] \textit{Analogy.}
The history space is the set of chemical elements
with their observed properties (atomic weight,
valence, melting point, oxidation states, reaction
products).  Mendeleev's 1869 table is the
explanation: each element is assigned a
(period, group) position based on atomic weight and
chemical-similarity grouping (later refined by atomic
number).  The descended map sends a property profile
to its table position; the target map predicts the
existence and properties of then-unknown elements
(gallium, germanium, scandium were subsequently
discovered with properties anticipated from table
gaps).  The compression is strict: the table position
summarises chemical-family behaviour without
recovering all numerical property values
\citep{Mendeleev1869Periodic}.

\item[\emph{Boltzmann {H}-theorem and
statistical-mechanical irreversibility (statistical
physics).}] \textit{Analogy.}
The history space is the microscopic phase-space
trajectory set of a dilute gas; the base records the
full microstate.  The explanation is the one-particle
distribution function \(f(\mathbf r,\mathbf v,t)\) and
the H-functional
\(H[f]=\int f\log f\,d^3v\,d^3r\).  Under the
molecular-chaos hypothesis, \(H\) is non-increasing
along Boltzmann-equation evolution; identifying
macroscopic entropy with \(S=-k_B H\), the target is
the macroscopic thermodynamic readout.  The
compression is strict: distinct microstates with the
same one-particle distribution share the same
macroscopic readout \citep{Boltzmann1872Weitere}.

\item[\emph{Hart's \emph{The Concept of Law} as
jurisprudential summary (jurisprudence).}] \textit{Analogy.}
The history space is the recorded legal practice of a
jurisdiction (statutes, decided cases, official
enactments); the base records the full legal corpus.
Hart's explanation distinguishes primary rules of
obligation from secondary rules of recognition,
change, and adjudication, together with the
identified rule of recognition.  The target, in
Hart's construction, is the legal-validity verdict on
candidate rules: a candidate counts as valid law iff
it satisfies the rule of recognition.  The target map
is exact because the rule of recognition is itself
the criterion of validity; the compression is strict
because many distinct legal corpora share the same
rule-system structure \citep{Hart1961Concept}.
\end{description}
\endgroup

\subsection*{\Fref{F31} Counterfactual Legitimacy}
\begingroup\small
\begin{description}
\item[\emph{Fogel counterfactual economic history of US
growth without railroads (cliometrics).}] \textit{Analogy.}
The context quotient collapses economic trajectories
to their factor-endowment class (capital, labour,
natural resources); the intervention abstraction
collapses transportation mixes to (railroad available,
wagon and canal substitute) categories; the route is
the social-savings simulation that re-prices freight
using non-rail substitutes for each route; the target
truth records whether the counterfactual 1890 GNP
exceeds a specified threshold (Fogel's calculation
gives a deviation of less than 5\% of actual).  The
legitimacy of the counterfactual is the \Fref{F31} statement:
the social-savings estimate descends through
factor-endowment and transportation-mix classes
\citep{Fogel1964RailroadsAmericanGrowth}.

\item[\emph{Legal ``but-for'' causation in tort law
(jurisprudence).}] \textit{Analogy.}
In common-law tort doctrine, but-for causation asks
whether the plaintiff's harm would have occurred but
for the defendant's act.  The context quotient
collapses factual histories to their legally cognisable
class; the intervention abstraction collapses
candidate acts to their legal characterisation;
admissibility encodes proximate-cause limits; the
route is the counterfactual reconstruction of ``what
would have happened'' without the defendant's act; the
target truth records whether the harm would still have
occurred.  The legitimacy of the but-for
counterfactual is the \Fref{F31} statement: the verdict
descends through legally cognisable
(context, act) classes, and overreaches (e.g.,
multiple-sufficient-cause cases where but-for fails)
trigger
doctrinal corrections such as the substantial-factor
test \citep{HartHonore1985Causation}.

\item[\emph{Off-policy evaluation in reinforcement
learning (machine learning).}] \textit{Analogy.}
An agent estimates the expected return of a target
policy \(\pi_e\) using trajectories collected under a
different behaviour policy \(\pi_b\).  The context
quotient collapses environments to their (state-action
transition kernel, reward function) class; the
intervention abstraction collapses target policies to
their induced visitation; admissibility requires the
support condition \(\pi_b>0\) wherever \(\pi_e>0\);
the route is the importance-sampling re-weighting; the
target truth records whether the estimated return of
\(\pi_e\) exceeds a specified threshold.  The
Precup-Sutton-Singh result gives conditions under
which the off-policy estimator descends from
behaviour-policy trajectories to a legitimate
counterfactual verdict on the target policy's value
\citep{PrecupSuttonSingh2000Eligibility}.

\item[\emph{Climate-model counterfactual emission
scenarios (climate science).}] \textit{Analogy.}
In atmospheric general-circulation modelling, the
context quotient collapses climate trajectories to
their baseline-forcing class (pre-industrial CO\(_2\),
solar and orbital parameters); the intervention
abstraction collapses scenarios to an idealised
greenhouse-gas forcing class (e.g., doubled
atmospheric CO\(_2\) concentration); the route is the
numerical integration of the coupled atmosphere
model under the chosen forcing; the target truth
records whether a specified climate diagnostic (e.g.,
equilibrium global-mean surface temperature change)
exceeds a threshold.  Manabe \& Wetherald 1967 is the
founding instance, with its calculation of equilibrium
climate sensitivity to doubled CO\(_2\) in a one-dimensional
radiative-convective model, a calculation the same authors repeated
in a general-circulation model in 1975; legitimacy
is the \Fref{F31} statement that the threshold-crossing
verdict descends through (baseline, idealised forcing)
classes \citep{ManabeWetherald1967Thermal,ManabeWetherald1975Doubling}.
\end{description}
\endgroup

\subsection*{\Fref{F33} Horizon Sufficiency / Holographic Compression}
\begingroup\small
\begin{description}
\item[\emph{{TCP/IP} encapsulation (network protocols).}] \textit{Analogy.}
A network message is encapsulated as a sequence of
nested headers wrapping a payload.  The interior is the
full packet state (headers and payload together); the
boundary readout records the IP and TCP header fields
(source/destination addresses, ports, sequence numbers,
flags, checksum).  The exterior classifier records the
IP protocol-number class (TCP, UDP, ICMP) and the
link-layer class.  Routing, flow-control, retransmission,
and demultiplexing decisions descend through the product
quotient: two packets sharing the same header fields and
the same protocol stack class receive the same routing
verdict, regardless of payload bytes
\citep{Postel1981RFC791,Postel1981RFC793}.

\item[\emph{Fisher sufficient statistic (mathematical
statistics).}] \textit{Analogy.}
For an exponential-family parametric model
\(p(x\mid\theta)\), a sufficient statistic
\(T(X_1,\dots,X_n)\) summarises a sample so that the
conditional distribution of the sample given \(T\) is
independent of \(\theta\).  The interior is the raw
sample space; the boundary readout is the sufficient
statistic; the exterior factor is a one-point set (an ancillary statistic is a function of the same sample and cannot fill \Fref{F33}'s separate exterior slot).  By the Fisher--Neyman factorization theorem,
likelihood-based inference functionals --- likelihood
ratio, score where defined, maximum-likelihood maximizing set (and any estimate selected solely from \(T\)), and
likelihood-equivalent functionals --- descend through \(T\), hence through
the product quotient \citep{Fisher1922Mathematical}.

\item[\emph{Cauchy's integral formula in complex analysis
(complex analysis).}] \textit{Analogy.}
For a function \(f\) holomorphic on a simply connected
domain \(D\subseteq\mathbb C\) and a point \(z_0\) in
the interior of a positively oriented simple closed
contour \(\gamma\subset D\),
\[
  f(z_0)=\frac{1}{2\pi i}\oint_\gamma \frac{f(z)}{z-z_0}\,
  dz.
\]
The interior parameterises the choice of contour
\(\gamma\) together with the holomorphic-function trace
\(f|_\gamma\) on it; the boundary readout records that
trace; the exterior classifier records the holomorphy
class of \(f\).  The target value \(f(z_0)\) at any
interior point descends exactly through the product
quotient: it is given by the Cauchy kernel integral
applied to the boundary trace
\citep{Ahlfors1979ComplexAnalysis}.
\end{description}
\endgroup

\subsection*{\Fref{F35} Scale Descent / Renormalization}
\begingroup\small
\begin{description}
\item[\emph{Multigrid methods on a nested mesh
hierarchy (numerical analysis).}] \textit{Analogy.}
For a linear elliptic PDE discretised on a nested
sequence of meshes
\(V_0\subset V_1\subset V_2\subset\cdots\) with mesh
sizes \(h_0>h_1>h_2>\cdots\), the carrier is the
finest-grid function space \(V_N\).  The scale quotient
is the composite restriction \(R_{n+1}\cdots R_N:V_N\to V_n\); the coarsening
is the one-step restriction \(R_n:V_n\to V_{n-1}\).  The discrete Galerkin variational
form \(A_n u_n=f_n\) at each level is the target law;
Galerkin coarsening \(A_{n-1}=R_n A_n P_n\) (with
prolongation \(P_n\) and \(R_n=P_n^{\!\top}\)) makes the coarse
operator the variational restriction of the fine one: for a symmetric
positive-definite \(A_n\) the coarse solution is the energy-norm
projection of the fine solution onto the coarse space, not its
restriction, so the coarse law is a descended approximation rather
than an exact descended law.  The multigrid V-cycle corrects it with
fine-level smoothing for asymptotically optimal solution costs \citep{Hackbusch1985Multigrid}.

\item[\emph{Coarse-grained molecular dynamics
(computational physics and materials).}] \textit{Analogy.}
For a biomolecular system, the carrier is the
all-atom phase-space trajectory.  The scale quotient
groups atoms into coarse-grained beads (e.g.,
MARTINI mapping of heavy atoms to lipid beads); the
coarsening iterates the grouping to even coarser
representations.  The target law at each level is a
structural-statistical observable matched between
the all-atom and coarse-grained models (radial
distribution functions, free-energy landscapes,
thermodynamic averages); the renormalised law is the
effective bead Hamiltonian or potential of mean force
chosen to reproduce those observables.  Scale descent
succeeds in the equilibrium structural regime and
becomes path-dependent for dynamical observables
sensitive to integrated-out fast atomic motions
\citep{Voth2008CoarseGraining}.

\item[\emph{Lucas critique on policy-invariant micro
parameters (macroeconomics).}] \textit{Analogy.}
The carrier is the set of individual agent decision
histories under a declared policy regime.  The scale
quotient collapses individual decisions to
macroeconomic aggregates (output, inflation,
employment); the coarsening sends aggregates under
one policy regime to aggregates under another via
macro-econometric extrapolation.  The Lucas critique
observes that reduced-form macro laws do not descend
through this coarsening when the policy change
itself alters agents' decision rules: the obstruction
set is non-empty.  Legitimate scale descent requires
identifying policy-invariant deep-structural micro
parameters and aggregating through them
\citep{Lucas1976Critique}.
\end{description}
\endgroup

\subsection*{\Fref{F36} Singular Limit Completion}
\begingroup\small
\begin{description}
\item[\emph{Prandtl boundary-layer matched asymptotic
expansion (fluid dynamics).}] \textit{Analogy.}
For the incompressible Navier-Stokes equations with
small viscosity \(\nu\), the set of limiting
processes is the family of solutions parameterised by
\(\nu\).  The limit quotient sends each solution to
its formal \(\nu\to 0\) limit; the target readout
records the velocity field.  The naive Euler limit
fails to satisfy the no-slip boundary condition;
Prandtl's boundary-layer theory introduces an inner
expansion in the rescaled wall-normal coordinate
\(y/\sqrt{\nu}\) and matches it to the outer Euler
solution.  The matched-asymptotic structure is the
canonical path-memory repair, carrying the
boundary-layer thickness as an additional quotient
parameter that completes the singular limit
\citep{Prandtl1904BoundaryLayer}.

\item[\emph{Critical-point scaling at second-order
phase transitions (statistical mechanics).}] \textit{Analogy.}
For a magnetic system near its critical temperature
\(T_c\), the set of limiting processes is the family
of approach paths to \((T,h)=(T_c,0)\) in the
temperature-field plane.  Order parameter,
susceptibility, and correlation length diverge or
vanish with universal exponents that depend on the approach path (\(|m|\propto(-t)^\beta\) on a selected ordered branch as \(t\uparrow0\) along \(h=0\), and \(|m|\propto|h|^{1/\delta}\) along \(t=0\)), with non-universal amplitudes.
The Widom scaling hypothesis is the path-memory
repair: it records the universal scaling functions
parameterised by \(h/|t|^{\beta\delta}\) (with
\(t=(T-T_c)/T_c\)), completing the critical-point
readout by labelling each approach direction
\citep{Stanley1971PhaseTransitions}.
\end{description}
\endgroup

\subsection*{\Fref{F37} Complementarity as Non-Joint Access}
\begingroup\small
\begin{description}
\item[\emph{Joint diagonalisation of non-commuting
Hermitians (linear algebra).}] \textit{Analogy.}
For a finite-dimensional Hilbert space \(H\) and two
Hermitian operators \(A,B\) on \(H\), the access
quotient \(q_A\) records the spectral labelling of
\(A\) in a basis and \(q_B\) records that of \(B\).
A joint admissible quotient is a single orthonormal
basis simultaneously diagonalising both operators.
By the spectral theorem, such a basis exists iff the
commutator \([A,B]=AB-BA\) vanishes.  When
\([A,B]\ne 0\), no admissible joint quotient exists;
the access quotients are \Fref{F37}-complementary within
the orthonormal-basis class
\citep{Halmos1974FiniteDimensional}.

\item[\emph{Intuitionistic vs classical logic
(mathematical logic).}] \textit{Analogy.}
For a propositional language \(L\) with atomic
propositions, the access quotient \(q_A\) sends a
formula \(\varphi\) to its intuitionistic-validity
class under the Brouwer-Heyting-Kolmogorov (BHK)
interpretation; \(q_B\) sends \(\varphi\) to its
classical-proof class requiring the law of excluded
middle (LEM) and proof by contradiction.  Within the
admissibility class of BHK-respecting proof systems,
no joint quotient validates LEM for undecidable
propositions \(P\): classical proof of
\(P\vee\neg P\) is admissible classically but lacks
a constructive witness.  The two proof-system
quotients are \Fref{F37}-complementary, formalised by
Heyting's intuitionistic logic
\citep{Heyting1930Logik}.

\item[\emph{{CAP} theorem in distributed systems
(distributed systems).}] \textit{Analogy.}
For a replicated distributed data store, \(q_A\)
records linearisable-consistency observations and
\(q_B\) records always-available reads and writes.
The admissibility constraint is partition tolerance:
the implementation must continue under arbitrary
network partitions.  Gilbert-Lynch's formalisation of
Brewer's CAP conjecture shows that, under network
partitions, no admissible implementation jointly
realises both \(q_A\) and \(q_B\); the access
quotients are \Fref{F37}-complementary within the
partition-tolerant implementation class
\citep{GilbertLynch2002CAP}.

\item[\emph{Russell's paradox under naive
comprehension (set-theoretic foundations).}] \textit{Analogy.}
For a candidate set \(R\), the access quotient
\(q_A\) imposes the self-membership predicate
\(x\in R\iff x\notin x\); the access quotient
\(q_B\) requires \(R\) to be a set in the admissible
class.  Substituting \(R\) for \(x\) yields
\(R\in R\iff R\notin R\), so no joint admissible
quotient exists under unrestricted comprehension.
The standard remedies (ZFC's separation schema, which forms \(\{x\in A:x\notin x\}\) only inside a given set \(A\), and Russell--Whitehead type
stratification) refine admissibility to block this
joint quotient \citep{Russell1902Letter}.
\end{description}
\endgroup

\subsection*{\Fref{F38} Arrow from Record Monotonicity}
\begingroup\small
\begin{description}
\item[\emph{Steno's law of superposition (geology).}] \textit{Analogy.}
The source history set is the configuration of
sediment-deposition processes; \(\kappa\) is the
geological-time evolution.  The record quotient sends
each deposition event to its stratum class (named bed,
geological epoch); the status readout is the
stratigraphic depth.  Steno's law of superposition
asserts that, in undisturbed sequences, older strata
lie below younger strata.  Folding, faulting, and
overturning are \Fref{F38}-obstructions that require
unfolding to recover the arrow \citep{Steno1669Prodromus}.

\item[\emph{Dendrochronological tree-ring sequence
(paleoclimate and archaeology).}] \textit{Analogy.}
The source history set is the sequence of annual
tree-ring formation events; \(\kappa\) records one
annual growth cycle.  The record quotient sends a
tree-ring to its ring-pattern signature; the status
readout is the calendrical year identified by master
tree-ring chronology cross-dating.  Each successive
ring is added outside the previous ones, producing a
strictly monotone chronological sequence.  Douglass's
master chronology cross-dates archaeological wood to
its year of growth via pattern matching against the
reference sequence \citep{Douglass1919ClimaticCycles}.

\item[\emph{Libby radiocarbon dating (archaeology and
nuclear chemistry).}] \textit{Analogy.}
The source history set is the moment of organism death
(after which exchange of \(^{14}\)C with the atmosphere
ceases); \(\kappa\) records radioactive decay over
time.  The record quotient sends a sample to its
measured \(^{14}\)C\(/^{12}\)C ratio; the status readout
is the radiocarbon age inferred from the exponential decay law;
conventional radiocarbon ages use the Libby half-life of 5568 years
(the measured value is about 5730 years) and are converted to
calendar ages by calibration against tree-ring and varve
chronologies.  The ratio decreases monotonically, providing a
chronological arrow up to about 50{,}000 years before present \citep{Libby1955RadiocarbonDating}.

\item[\emph{Historical phonological reconstruction.}] \textit{Analogy.}
The source history set is the lexicon of a
reconstructed proto-language; the target is the
lexicon of an attested daughter language; \(\kappa\)
records the regular sound-change rules carrying
proto-forms forward to daughter reflexes.  The record
quotient sends a proto-form to its phonological class;
the status readout is the chronological layer in the
family tree, increasing from the proto-language toward
each attested daughter.  Brugmann's regularity
principle asserts that sound changes are exceptionless
within a defined phonological environment, providing a
monotone forward arrow along the family tree
\citep{Brugmann1886Grundriss}.
\end{description}
\endgroup

\subsection*{\Fref{F39} Entropy as Fiber Volume}
\begingroup\small
\begin{description}
\item[\emph{Herfindahl-Hirschman market-concentration
index (industrial economics).}] \textit{Analogy.}
For a market with firms having shares \(s_i\)
(\(\sum_i s_i=1\)), the access quotient sends each
customer transaction to the firm serving it.  The
Herfindahl-Hirschman index
\(\mathrm{HHI}=\sum_i s_i^2\) is the probability that
two random transactions are served by the same firm.
Under firm grouping (mergers, parent-company
aggregation), nonnegativity of the cross terms in the square of a grouped share
forces the sum of squared shares to non-decrease, so HHI is monotone under
coarsening.  The US Department of Justice merger
guidelines use HHI thresholds to identify market
concentration concerns \citep{Hirschman1964HHI}.

\item[\emph{Logarithm of the multinomial coefficient
(combinatorics).}] \textit{Analogy.}
For \(N\) distinguishable items distributed into \(k\)
labelled cells, the fine quotient sends an assignment
to its cell-content profile
\((n_1,\dots,n_k)\); the fine fiber has cardinality
\(\binom{N}{n_1,\dots,n_k}=N!/\prod_i n_i!\).
Coarsening merges two cells \(i,j\) into a single
labelled cell of size \(n_{\{i,j\}}=n_i+n_j\); the
coarse fiber is the disjoint union of fine fibers
with \(n_i'+n_j'=n_{\{i,j\}}\), so its cardinality is
\(\sum_{n_i'+n_j'=n_{\{i,j\}}}
\binom{N}{n_i',n_j',\dots}\), at least any single
summand.  The combinatorial fiber-volume entropy is
thus monotone non-decreasing under cell coarsening
\citep{Knuth1997TAOCP}.

\item[\emph{Kingman coalescent effective population
size (population genetics).}] \textit{Analogy.}
For a haploid neutral genealogical model, the access
quotient sends a genealogy of \(n\) sampled lineages
to its lineage-merger pattern.  The Kingman coalescent
gives pairwise coalescence rate \(1/N_e\), where
\(N_e\) is the effective population size; with \(n\) lineages the expected waiting time to the next merger is \(N_e/\binom n2\).  Changing \(N_e\) rescales coalescent time rather than coarsening a quotient, so this is only a loose analogue of the \Fref{F39} inequality; coarsening a finite carrier of merger patterns with fiber-volume entropy gives an instance \citep{Kingman1982Coalescent}.
\end{description}
\endgroup

\subsection*{\Fref{F42} Irreducibility / No Short Closed Description}
\begingroup\small
\begin{description}
\item[\emph{Quine-McCluskey minimum sum-of-products
form (digital-logic design).}] \textit{Analogy.}
For a Boolean function \(f:\{0,1\}^n\to\{0,1\}\), the
admissible compression class is sum-of-products
expressions; routes are Boolean-algebra equivalences
(De Morgan, absorption, distribution).  Quine-McCluskey
yields a minimum-literal expression by generating all prime implicants and then solving a minimum-weight cover problem, weighting each prime implicant by its number of literals (NP-hard in general) for the minterms not covered by essential prime implicants; the
function is \Fref{F42}-irreducible at the minimum-cover
expression \citep{Quine1952Boolean}.

\item[\emph{Information-theoretic
\(\Omega(n\log n)\) sorting lower bound (analysis of
algorithms).}] \textit{Analogy.}
For comparison-sorting on \(n\) distinct keys, the
admissible compression class is comparison-based
decision trees; routes are two-key comparisons; the
sorted permutation is the target.  A correct decision
tree must have at least \(n!\) leaves, so its
worst-case depth is \(\Omega(\log n!) = \Omega(n\log
n)\).  No asymptotically shallower comparison-based
decision tree exists; the \Fref{F42} irreducibility is the
unconditional \(\Omega(n\log n)\) lower bound
\citep{Knuth1998TAOCPv3}.

\item[\emph{Minimal genome for self-replicating cell
(synthetic biology).}] \textit{Analogy.}
Hutchison and colleagues designed and synthesized
JCVI-syn3.0, a \emph{Mycoplasma} genome of 531
kilobases encoding 473 genes that supports
self-replication under declared laboratory growth
conditions.  The design retained genes classified as
essential and also quasi-essential genes needed for
robust growth \citep{Hutchison2016MinimalGenome}.
The analogy is with the search for a short closed
description within a declared class of designs.  The
experiment tests only part of that class, so it does
not establish the minimality that \Fref{F42}
requires.

\item[\emph{Crystallographic irreducible Brillouin
zone (solid-state physics).}] \textit{Analogy.}
For a periodic crystalline solid with space group
\(G\), the admissible compression is the
restriction of crystal momentum \(\mathbf k\) to the
irreducible Brillouin zone (IBZ), the fundamental
domain of the reciprocal-lattice point-group action;
routes are point-group operations.  The
Bouckaert-Smoluchowski-Wigner analysis shows the
IBZ is the minimal subset of the full Brillouin
zone that, together with symmetry operations,
generates all distinct electronic states.  The IBZ
is \Fref{F42}-irreducible by construction
\citep{BouckaertSmoluchowskiWigner1936Theory}.
\end{description}
\endgroup

\subsection*{\Fref{F43} Continuum Emergence}
\begingroup\small
\begin{description}
\item[\emph{p-adic completion in number theory
(algebraic number theory).}] \textit{Analogy.}
For a prime \(p\), the discrete approximations are
the finite quotient rings \(\mathbb Z/p^i\mathbb Z\)
with reduction maps
\(\phi_{ij}:\mathbb Z/p^j\to\mathbb Z/p^i\).  The
candidate limit is the p-adic integers
\(\mathbb Z_p=\varprojlim\mathbb Z/p^i\mathbb Z\),
with projections to each finite quotient.  Tower
coherence is congruence compatibility, cofinality
is directedness of p-power indices, and vanishing
residuals correspond to convergence in the p-adic
metric.  \(\mathbb Z_p\) is the p-adic completion of
\(\mathbb Z\); the non-archimedean continuum
\(\mathbb Q_p=\mathbb Z_p[1/p]\) is the p-adic completion of
\(\mathbb Q\), distinct from the archimedean reals
\citep{Hensel1908Theorie}.

\item[\emph{Thermodynamic limit in statistical
mechanics (statistical physics).}] \textit{Analogy.}
For a classical-spin lattice model on a finite
volume \(\Lambda\subset\mathbb Z^d\), the discrete
approximations are the finite-volume Gibbs kernels at inverse
temperature \(\beta\), which condition on the configuration outside
the volume and are consistent under nesting of volumes; finite-volume
Gibbs measures with one fixed boundary condition are not consistent
under restriction.  The candidate limit is
the set of Dobrushin-Lanford-Ruelle (DLR) states
on \(\{+1,-1\}^{\mathbb Z^d}\), with projections
to finite-volume marginals.  Tower coherence is
the DLR consistency condition, cofinality is
exhaustion by finite volumes, and vanishing
residuals correspond to absence of unstable
boundary contributions.  The set of infinite-volume
Gibbs states emerges as the \Fref{F43} stable limit quotient; it need
not be a singleton, and at low temperature in dimension \(d\ge 2\)
the Ising model has several (phase coexistence)
\citep{Ruelle1969StatisticalMechanics}.

\item[\emph{Krull absolute Galois group as inverse
limit (field theory).}] \textit{Analogy.}
For a field \(K\) with separable closure
\(K^{\mathrm{sep}}\), the discrete approximations
are finite Galois groups
\(\mathrm{Gal}(L_i/K)\) for finite Galois
extensions \(L_i\), with restriction maps for
\(L_i\subseteq L_j\).  The candidate limit is the
absolute Galois group
\(\mathrm{Gal}(K^{\mathrm{sep}}/K)
=\varprojlim_i\mathrm{Gal}(L_i/K)\), with
projections by restriction.  Tower coherence is
compatibility of restrictions, cofinality is
directedness of finite Galois extensions, and
vanishing residuals correspond to profinite
separation by the finite quotients.  Krull's
infinite Galois theory establishes the absolute
Galois group with its profinite (Krull) topology
as the \Fref{F43} stable limit quotient
\citep{Krull1928Galois}.

\item[\emph{Born-Huang continuum elasticity from
atomic lattice (solid-state physics).}] \textit{Analogy.}
For a crystalline solid, the discrete
approximations are atomic-lattice models at
finer-and-finer lattice spacings; comparison maps
are inter-lattice coarse-graining operations.
The candidate limit is the continuum elasticity
theory governing displacement fields on
\(\mathbb R^3\); projections restrict continuum
solutions to lattice points.  Tower coherence is
compatibility of inter-lattice strain
representations, cofinality is the existence of
sufficiently fine lattices for any continuum
scale, and vanishing residuals correspond to the
long-wavelength approximation being valid.  The
Born-Huang derivation of elastic constants from
interatomic potentials is the canonical
solid-state \Fref{F43} continuum emergence
\citep{BornHuang1954Dynamical}.
\end{description}
\endgroup

\subsection*{\Fref{F44} Criticality / Phase Boundary}
\begingroup\small
\begin{description}
\item[\emph{Reynolds laminar-turbulent transition
(fluid dynamics).}] \textit{Analogy.}
For incompressible smooth-pipe flow, the control
space is the Reynolds number
\(\mathrm{Re}=\rho v D/\mu\); the phase readout
records laminar vs turbulent.  Pipe-flow
experiments since Reynolds's identify a critical
\(\mathrm{Re}_c\approx 2300\) as the phase boundary at which laminar
flow loses empirical robustness to finite-amplitude perturbations at
ordinary inflow disturbance levels; with carefully quieted inflow,
laminar flow persists far above this value, to about 13{,}000 in
Reynolds's own experiments.  The
empirical robustness margin vanishes, the
intermittent-puff response is non-uniform, and the
residual status is described by transition
intermittency statistics
\citep{Reynolds1883Turbulent}.
\end{description}
\endgroup

\subsection*{\Fref{F47} Fine-Tuning as Small Selector Region}
\begingroup\small
\begin{description}
\item[\emph{Pharmaceutical therapeutic window
(clinical pharmacology).}] \textit{Analogy.}
The moduli space is the set of dose-and-timing
protocols; the readout is the steady-state plasma
concentration; the target predicate restricts to
concentrations between the minimum effective
concentration (MEC) and the maximum tolerated
concentration (MTC).  The therapeutic-window selector
region is small for drugs with narrow therapeutic
indices (warfarin, lithium, digoxin); realized
fine-tuning is the clinically attained dose protocol
placing plasma concentration within the window
\citep{GoodmanGilman2018Pharmacology}.

\item[\emph{Mechanical tolerance stack-up gap
(manufacturing engineering).}] \textit{Analogy.}
The moduli space is the set of component dimension
tuples within their tolerance bands; the readout is
the resulting critical assembly dimension; the
target predicate restricts to assemblies meeting the
critical-dimension specification.  The selector
region shrinks as the assembly chain lengthens;
realized fine-tuning is the actually manufactured
component-dimension tuple landing inside the
allowable region \citep{Bjorke1989Tolerance}.

\item[\emph{Hutchinson ecological niche (ecology).}] \textit{Analogy.}
The moduli space is the multi-dimensional
environmental parameter space (temperature,
humidity, pH, salinity, prey abundance); the readout
is the species' population growth rate; the target
predicate restricts to non-negative growth.  Within
the declared environmental envelope for the species,
the niche occupies a small fraction of total moduli
volume; realized fine-tuning is the species' actual
occupancy of the niche
\citep{Hutchinson1957Concluding}.

\item[\emph{Nyquist closed-loop stability region
(control theory).}] \textit{Analogy.}
The moduli space is the set of (plant,
controller-parameter) combinations; the readout is
the closed-loop pole location in the complex plane;
the target predicate restricts to all poles having
negative real part.  The Nyquist criterion
characterises the stability selector region: the
closed-loop system is stable iff the number of
counter-clockwise encirclements of \(-1+0i\) equals
the number of open-loop right-half-plane poles.  For
systems with large gain, delay, or open-loop
instability, the stability region can be small,
instantiating \Fref{F47} fine-tuning of controller
parameters \citep{Nyquist1932Regeneration}.
\end{description}
\endgroup

\subsection*{\Fref{F49} Common-Source / Nonlocal Correlation Normal Form}
\begingroup\small
\begin{description}
\item[\emph{Doolittle lateral gene transfer
(microbial phylogenetics).}] \textit{Analogy.}
For prokaryotic phylogenomics, the source quotient
sends a gene-distribution configuration to its
vertical-inheritance most-recent-common-ancestor
class; the readouts are per-gene tree topologies
and presence patterns.  Vertical inheritance gives
CommonSource correlations across co-inherited
genes; Doolittle observed that horizontal gene
transfer events mediated by phages, plasmids, and
conjugation break the common-source pattern,
constituting an interface channel.  A pattern is LocallyExplainable when it factors through vertical descent, a transfer interface, or a direct route; \Fref{F49} gives this as a disjunction, not an additive decomposition \citep{Doolittle1999Phylogenetic}.

\item[\emph{Plagiarism detection via shared-source
fingerprinting (digital forensics).}] \textit{Analogy.}
For pairs of documents, the source quotient sends
a pair to its shared-source text together with
declared edit / paraphrase / template
transformations; each document is fingerprinted via
n-gram or hashed-shingle representations; the
correlation readout is an overlap metric (Jaccard
similarity, MinHash).  CommonSource holds when the
overlap descends through the
shared-source-plus-declared-transformations
quotient; the interface
captures legitimate citation, common literature,
and formulaic templates not classified as
plagiarism \citep{Heintze1996Scalable}.

\item[\emph{Mosteller-Wallace stylometric authorship
attribution (statistical humanities).}] \textit{Analogy.}
For the disputed Federalist Papers, the source
quotient assigns a paper to a candidate author; the
readouts are author-conditioned function-word
frequency distributions
\(P(\mathrm{wordfreq}\mid\sigma)\) estimated from
each candidate's known corpus.  The Bayesian
posterior \(P(\sigma\mid\mathrm{wordfreq})\) is a probabilistic analogue of CommonSource, which \Fref{F49} covers only as a deterministic factorization; joint
authorship or formulaic shared templates
instantiate the interface, and residual anomalies
index unmodeled influences
\citep{MostellerWallace1964Federalist}.

\item[\emph{Common-cause failure analysis in
fault trees (safety engineering).}] \textit{Analogy.}
For a safety-critical system, the source quotient
sends a failure event to its underlying
common-cause basic event (loss of off-site power,
fire,
earthquake).  A joint failure fixed by a shared basic event is the deterministic CommonSource pattern; statistical common-cause correlation lies outside \Fref{F49}.  Fault-tree analysis
identifies the basic events that constitute the
source, codifying methodology for nuclear-plant
reliability and risk assessment
\citep{Vesely1981FaultTree}.
\end{description}
\endgroup

\subsection*{\Fref{F50} Vacuum / Background Selection}
\begingroup\small
\begin{description}
\item[\emph{Ostwald rule of stages: polymorph
selection (crystallography).}] \textit{Analogy.}
For a supersaturated solution, the moduli quotient
collapses crystallisation trajectories sharing the
same crystalline polymorph; the descended
background records the polymorph type.  Ostwald's
rule of stages predicts kinetic selection of a
metastable polymorph close in free energy to the
supersaturated state.  The vacuum predicate
identifies the ultimately stable polymorph as the
singleton target-formed modulus
\citep{Ostwald1897Studien}.

\item[\emph{Zermelo axiomatisation (foundations
of mathematics).}] \textit{Analogy.}
The closure space is the set of formal theories
purporting to be foundations of mathematics; the
moduli quotient sends a theory to its axiom-system
class (Zermelo 1908 set theory,
Fraenkel-Skolem-extended ZF, ZFC, NBG, intuitionistic
ZF,
Martin-Löf type theory).  Zermelo's 1908
axiomatisation selected a specific moduli value —
Zermelo set theory with separation and choice —
as a singleton background; subsequent extensions
(Fraenkel Replacement, Foundation, large cardinals)
refine that selection within the same \Fref{F50}
background-moduli framework
\citep{Zermelo1908Untersuchungen}.

\item[\emph{Hund's rules for atomic ground state
(atomic physics).}] \textit{Analogy.}
For a multi-electron atom, the moduli quotient
sends Pauli-allowed configurations to their
\((S,L,J)\) term-symbol class; the descended
background is the ground-state term symbol.
Hund's rules select the ground state: (i) maximize
total \(S\); (ii) for given \(S\), maximize \(L\);
(iii) for given \((S,L)\), minimize \(J\) if the
subshell is less than half-filled, otherwise
maximize \(J\).  Within the LS-coupling regime, the
vacuum predicate is the Hund-selected ground term
symbol \citep{Hund1925Linienspektren}.

\item[\emph{Pittendrigh circadian phase entrainment
(chronobiology).}] \textit{Analogy.}
For a circadian oscillator under a periodic
zeitgeber, the moduli quotient sends oscillator
trajectories sharing the same entrained phase to a
class; the descended background records the
entrained phase.  Pittendrigh's phase-response-curve
analysis explains how zeitgebers select a definite
entrained phase from the continuous phase moduli
space; the vacuum predicate is the biologically
realised entrained phase
\citep{Pittendrigh1960Circadian}.
\end{description}
\endgroup

\subsection*{\Fref{F51} Unification as Common Refinement}
\begingroup\small
\begin{description}
\item[\emph{Cartesian analytic geometry
(mathematics).}] \textit{Analogy.}
The parent closure is Cartesian analytic geometry;
the children are Euclidean geometry (synthetic
proofs) and algebra (symbolic equations).  Cartesian
coordinates supply interpretations in both
directions: every coordinate-representable
algebraic curve maps to a geometric object and
vice versa.  Euclidean theorems become algebraic
identities on coordinates, and algebraic equations
acquire geometric interpretation as curves and
surfaces \citep{Descartes1637Geometrie}.

\item[\emph{Banach functional analysis (functional
analysis).}] \textit{Analogy.}
The parent closure is Banach's framework of complete
normed vector spaces with bounded linear operators;
the children are finite-dimensional linear algebra
and classical function-space analysis.  Both embed
as concrete Banach spaces, and theorems
(Hahn-Banach, open mapping, closed graph, uniform
boundedness) supply unified results across both
children within the framework's declared scope
\citep{Banach1932Theorie}.

\item[\emph{L\'evi-Strauss structural anthropology
(anthropology).}] \textit{Analogy.}
The parent closure is L\'evi-Strauss's structural
analysis of myth via mythemes (binary oppositions,
mediating figures); the children are
ethnographically distinct mythological corpora.
Each corpus is interpreted in terms of its
structural elements, and the parent identifies
mythematic isomorphisms across geographically
isolated cultures \citep{LeviStrauss1958Anthropologie}.

\item[\emph{Huxley modern synthesis (evolutionary
biology).}] \textit{Analogy.}
The parent closure is the modern synthesis (Huxley
1942), unifying Mendelian genetics and Darwinian
natural selection together with cytology,
paleontology, and biogeography.  Population
genetics supplies the common refinement:
allele-frequency dynamics under Mendelian
inheritance extend, with selection coefficients, to a
unified
predictive framework recovering both Mendelian
ratios and Darwinian adaptive change as special
cases \citep{Huxley1942Evolution}.
\end{description}
\endgroup

\subsection*{\Fref{F52} Closure Completeness Boundary}
\begingroup\small
\begin{description}
\item[\emph{{PMBOK} project-scope management
(project management).}] \textit{Analogy.}
The residuals are deliverables and acceptance
criteria; directions are project activities and
stakeholder requests.  The PMBOK scope statement
declares deliverables (\emph{inScope}), exclusions
(\emph{namedOutside}), assumptions, constraints,
and acceptance criteria; scope creep is a
stakeholder request that is neither.  The
meta-audit is project-sponsor sign-off
\citep{PMI2017PMBOK}.

\item[\emph{International Building Code
jurisdictional scope (construction codes).}] \textit{Analogy.}
The residuals are code-mandated safety,
structural, fire-protection, and accessibility
requirements; directions are declarable
construction features.  The IBC declares scope by
occupancy classification, construction type, and
jurisdictional adoption, with detached one- and two-family dwellings and townhouses, each not more than three stories above grade plane and, for townhouses, with a separate means of egress, referred to the International Residential Code (IBC 101.2) \emph{namedOutside}.  The meta-audit is the
building inspector's certificate of occupancy
\citep{ICC2021IBC}.

\item[\emph{Patent-claim scope under {M}arkman
hearings (intellectual-property law).}] \textit{Analogy.}
The residuals are patent-claim terms requiring
construction; directions are declarable claim
scopes and accused infringing embodiments.  Under
\emph{Markman v. Westview Instruments}, claim
construction is a matter of law: each disputed term
is construed (\emph{inScope}) or held indefinite
under \S112 (\emph{namedOutside} the valid claim).
The meta-audit is the Federal Circuit affirmance of
district-court claim construction
\citep{Markman1996Westview}.

\item[\emph{{FDA} drug-label indication scope
(regulatory affairs).}] \textit{Analogy.}
The residuals are safety and efficacy items required
on the FDA-approved label; directions are medical
uses.  21 CFR 201.57 prescribes the format and
content of the ``Indications and Usage'' section,
with FDA-approved indications \emph{inScope} and
off-label uses \emph{namedOutside}.  The meta-audit
is FDA approval of the (Supplemental) New Drug
Application \citep{FDA21CFR20157}.
\end{description}
\endgroup

\subsection*{Direct CER no-go (\Cref{thm:CER-direct-no-go})}
\begingroup\small
\begin{description}
\item[\emph{Preclinical reproducibility record
(research methodology).}] \textit{Analogy.}
The carrier \(C\) is the reported experimental
result; \(\mathfrak S\) imposes the
reproducibility audit (lab notebook, raw data,
randomisation log, blinded scoring).  The
closure-relevant datum \(d\) is the
replication-enabling record.  Begley and Ellis
report that Amgen scientists confirmed only 6 of 53 ``landmark''
preclinical findings, and attribute the failures in part to missing
controls, unblinded scoring and selective reporting.  Read through
the theorem, the bench experiments were bare-realised yet the record
required to close the published claim was inadequate, so the
reduct satisfies \(\operatorname{BareReal}
_{\mathfrak S}\) but the reproducibility
residual has no satisfying record
\citep{Begley2012Drug}.

\item[\emph{Software bill of materials (software
supply chain).}] \textit{Analogy.}
The carrier \(C\) is the compiled binary in
production; \(\mathfrak S\) imposes the
supply-chain audit (component inventory, licence
terms, vulnerability mapping, build provenance).
The closure-relevant datum \(d\) is the SBOM
listing every linked component with supplier, component name, version, other unique identifiers and dependency relationships, per the NTIA minimum elements.
The binary runs without the SBOM, so the reduct
satisfies \(\operatorname{BareReal}_{\mathfrak S}\)
yet the retained claim ``free of
known-vulnerable components'' has a residual
whose audit datum (the SBOM) was never produced
\citep{NTIA2021SBOM}.

\item[\emph{Photometric flux without calibration
log (observational astronomy).}] \textit{Analogy.}
The carrier \(C\) is the recorded photon count on
a CCD; \(\mathfrak S\) imposes the photometric
audit (flat field, dark current, atmospheric
extinction, zero-point).  The closure-relevant
datum \(d\) is the per-exposure calibration log
binding ADU counts to absolute flux, as in the
SDSS \"ubercal solution.  The raw integration is
bare-realised even when \(d\) is absent, so the
reduct satisfies \(\operatorname{BareReal}
_{\mathfrak S}\) yet the retained claim ``object
has AB magnitude \(m\)'' has a residual whose
audit datum is the missing calibration
\citep{Padmanabhan2008Improved}.

\item[\emph{Chain of title for real property
(real-property law).}] \textit{Analogy.}
The carrier \(C\) is the parcel together with
the occupant's actual possession; \(\mathfrak S\)
imposes the marketable-title audit (recorded
deeds, mortgages, easements, liens).  The
closure-relevant datum \(d\) is the unbroken
chain of recorded conveyances from a sovereign
source.  Adverse occupancy bare-realises the
parcel even when an unrecorded prior deed or
undischarged mortgage breaks the chain, so the
reduct satisfies \(\operatorname{BareReal}
_{\mathfrak S}\) but the retained claim
``marketable title'' has a residual whose
required record is the missing deed
\citep{Palomar2003Patton}.
\end{description}
\endgroup

\subsection*{Converse CER no-go (\Cref{thm:CER-converse-no-go})}
\begingroup\small
\begin{description}
\item[\emph{Ratified RFC without deployed
implementation (Internet standards).}] \textit{Analogy.}
The carrier \(C\) is the standards-track RFC
considered as a closure-of-record: working-group
consensus, IETF Last Call, IESG approval, RFC
Editor publication, IANA registry allocation per
RFC 2026.  Under \(\mathfrak S\) the standards
process is closed at publication.  When no
production implementation deploys the protocol
(orphan or stillborn RFCs), no Internet host
realises \(C\), so \(C=\mathrm{Cl}_{\mathfrak S}
(C)\) yet \(\neg\operatorname{BareReal}_{\mathfrak
S}(C)\) \citep{Bradner1996RFC2026}.

\item[\emph{Statute fallen into desuetude
(statutory law).}] \textit{Analogy.}
The carrier \(C\) is the enacted statute
considered as a closure-of-record: bill passed by
both chambers, presented and signed (or enacted
over veto), codified in the official compilation.
Under \(\mathfrak S\) the legislative procedure is
closed at codification.  When prosecutors decline
to charge, courts decline to apply, and citizens
disregard the statute over sustained time, no
enforcement instance realises \(C\), so
\(C=\mathrm{Cl}_{\mathfrak S}(C)\) yet \(\neg
\operatorname{BareReal}_{\mathfrak S}(C)\); desuetude is not a recognised doctrine in
most United States jurisdictions (West Virginia is the usual
exception), and Calabresi proposes that courts be allowed to treat
such obsolete statutes as open to judicial revision \citep{Calabresi1982Common}.

\item[\emph{Superconducting Super Collider cancellation.}] \textit{Analogy.}
The carrier \(C\) is the proposed 40\,TeV
(centre-of-mass) proton--proton collider considered as a
closure-of-record: Conceptual Design Report,
Site-Specific
Conceptual Design Report, technical design
reports, DOE design reviews, environmental impact
statement, congressional authorisations.  Under
\(\mathfrak S\) the megaproject design-review
audit is closed at the review-completion
boundary.  Congressional cancellation in 1993
terminated construction before operations, so no
running collider instantiates \(C\).  Hence
\(C=\mathrm{Cl}_{\mathfrak S}(C)\) yet
\(\neg\operatorname{BareReal}_{\mathfrak S}(C)\)
\citep{Riordan2015Tunnel}.

\item[\emph{Standard {ML} signature without a
structure (programming-language type systems).}] \textit{Analogy.}
The carrier \(C\) is a Standard ML signature
considered as a closure-of-record: abstract type
declarations, value specifications with types,
exception declarations, all type-checked against
the SML definition.  Under \(\mathfrak S\) the
type-system audit is closed at signature
elaboration.  Without a matching structure binding
no runtime value inhabits the abstract types and
no function body computes the declared values, so
\(C=\mathrm{Cl}_{\mathfrak S}(C)\) yet \(\neg
\operatorname{BareReal}_{\mathfrak S}(C)\)
\citep{MacQueen1984Modules}.
\end{description}
\endgroup

\subsection*{\Fref{F53} Closure--Reality Boundary resolution}
\begingroup\small
\begin{description}
\item[\emph{Audited financial statements (financial reporting).}]
\textit{Analogy.} Closing a ledger's inference rules does not supply missing
transaction evidence or an auditor's attestation. Accounting needs actual
witnesses satisfying the declared obligations; acceptance of an opinion is
an additional consumer requirement \citep{FASB2009ASC}.

\item[\emph{ACM Artifact Review and Badging (computational reproducibility).}]
\textit{Analogy.} Deriving consequences of an artifact manifest does not
create a missing executable or reproduction log. Those witnesses must be
supplied and checked before the relevant policy can award a badge
\citep{ACM2020Artifact}.

\item[\emph{CompCert formally verified compiler (programming-language verification).}]
\textit{Analogy.} Inference closure under proof rules does not create a
missing semantic-preservation certificate. The actual proof must satisfy
the original source-to-target specification \citep{Leroy2009Formal}.

\item[\emph{End-to-end verifiable elections (election integrity).}]
\textit{Analogy.} Closing the declared public record under inference does
not create missing shuffle or decryption proofs. Actual supplied certificates
and the verification policy determine whether the original tally obligation
is delivered \citep{Adida2008Helios}.
\end{description}
\endgroup
